\documentclass[11pt,a4paper]{article}

% shared/macros.tex
% Shared notation for all surface-partition math documents.
% Include with: % shared/macros.tex
% Shared notation for all surface-partition math documents.
% Include with: % shared/macros.tex
% Shared notation for all surface-partition math documents.
% Include with: \input{../shared/macros}

% -------------------------------------------------------------------------
% Packages (include here so each document only needs \input{../shared/macros})
% -------------------------------------------------------------------------
\usepackage{amsmath,amssymb,amsthm}
\usepackage{bm}           % \bm for bold math
\usepackage{mathtools}    % \coloneqq, dcases, etc.
\usepackage{booktabs}     % \toprule etc. in tables
\usepackage{hyperref}
\usepackage{geometry}
\usepackage{microtype}
\usepackage{parskip}
\usepackage{xcolor}
\usepackage{tcolorbox}
\usepackage{enumitem}
\usepackage{float}             % [H] float placement

\geometry{a4paper, margin=2.5cm}

% -------------------------------------------------------------------------
% Theorem-like environments
% -------------------------------------------------------------------------
\theoremstyle{definition}
\newtheorem{defn}{Definition}[section]
\newtheorem{remark}{Remark}[section]

\theoremstyle{plain}
\newtheorem{prop}{Proposition}[section]

% -------------------------------------------------------------------------
% Mesh and geometry
% -------------------------------------------------------------------------
\newcommand{\V}{\mathbf{V}}            % vertex array V[i] in R^dim
\newcommand{\F}{\mathbf{F}}            % face (triangle) array
\newcommand{\vone}[1]{v_{1,#1}}       % first vertex index of VP i
\newcommand{\vtwo}[1]{v_{2,#1}}       % second vertex index of VP i

% -------------------------------------------------------------------------
% Variable points
% -------------------------------------------------------------------------
\newcommand{\lam}{\lambda}             % scalar VP parameter
\newcommand{\Lam}{\boldsymbol{\lambda}}% full parameter vector
\newcommand{\vp}[1]{\mathbf{p}_{#1}}  % 3D position of VP i
\newcommand{\ed}[1]{\mathbf{d}_{#1}}  % edge direction d_i = V[v1_i] - V[v2_i]

% -------------------------------------------------------------------------
% Steiner / triple-point
% -------------------------------------------------------------------------
\newcommand{\Sv}{\mathbf{S}}           % Steiner (Fermat-Torricelli) point
\newcommand{\Ki}{K_{i}}               % projection matrix (I-n_i n_i^T)/r_i
\newcommand{\Mmat}{M}                 % sum of K_i matrices at a triple point

% -------------------------------------------------------------------------
% Normals and unit vectors
% -------------------------------------------------------------------------
\newcommand{\nhat}{\hat{\mathbf{n}}}  % unit normal to a triangle
\newcommand{\Dhat}{\hat{\boldsymbol{\Delta}}}  % unit segment direction

% -------------------------------------------------------------------------
% Operations
% -------------------------------------------------------------------------
\newcommand{\norm}[1]{\left\|#1\right\|}
\newcommand{\abs}[1]{\left|#1\right|}
\newcommand{\ip}[2]{\langle #1,\, #2 \rangle}

% Projection perpendicular to a unit vector \uhat
\newcommand{\Proj}[1]{\bigl(\mathbf{I} - #1 #1^\top\bigr)}

% argmax operator (winner-take-all assignment in Phase 1 documents)
\DeclareMathOperator*{\argmax}{arg\,max}
\DeclareMathOperator*{\argmin}{arg\,min}

% -------------------------------------------------------------------------
% Calculus shorthands
% -------------------------------------------------------------------------
\newcommand{\pd}[2]{\dfrac{\partial #1}{\partial #2}}
\newcommand{\pdd}[3]{\dfrac{\partial^2 #1}{\partial #2\,\partial #3}}
\newcommand{\grad}{\nabla}
\newcommand{\Hess}[1]{\nabla^2 #1}

% -------------------------------------------------------------------------
% Optimization problem objects
% -------------------------------------------------------------------------
\newcommand{\Preg}{P_{\mathrm{reg}}}  % regular perimeter
\newcommand{\Pst}{P_{\mathrm{st}}}   % Steiner perimeter
\newcommand{\Areg}{A^{\mathrm{reg}}} % regular cell area
\newcommand{\Ast}{A^{\mathrm{st}}}   % Steiner area contribution
\newcommand{\Atgt}{\bar{A}}           % target area per cell
\newcommand{\Lagr}{\mathcal{L}}       % Lagrangian
\newcommand{\objfac}{\sigma}          % IPOPT objective scaling factor

% -------------------------------------------------------------------------
% Code cross-references (for "In the code..." callouts)
% -------------------------------------------------------------------------
\newcommand{\code}[1]{\texttt{#1}}
\newcommand{\codefile}[1]{\texttt{\small #1}}

% -------------------------------------------------------------------------
% Threshold dynamics / approach B (10-mbo-auction-dynamics)
% -------------------------------------------------------------------------
\newcommand{\Dlump}{D}                 % lumped mass matrix diag(v)
\newcommand{\Atau}{A_\tau}             % discrete diffusion operator (M + tau K)^{-1} M
\newcommand{\Etau}{E_\tau}             % heat-content (threshold-dynamics) energy
\newcommand{\Rcell}{R_{\mathrm{cell}}} % radius of the equal-area disc sqrt(Abar/pi)
\newcommand{\hmean}{h_{\mathrm{mean}}} % mean edge length
\newcommand{\hmax}{h_{\mathrm{max}}}   % max edge length
\newcommand{\Gt}{G_t}                  % heat kernel at time t
\newcommand{\Ktau}{K^{\mathrm{res}}_\tau} % resolvent (backward-Euler) kernel


% -------------------------------------------------------------------------
% Packages (include here so each document only needs % shared/macros.tex
% Shared notation for all surface-partition math documents.
% Include with: \input{../shared/macros}

% -------------------------------------------------------------------------
% Packages (include here so each document only needs \input{../shared/macros})
% -------------------------------------------------------------------------
\usepackage{amsmath,amssymb,amsthm}
\usepackage{bm}           % \bm for bold math
\usepackage{mathtools}    % \coloneqq, dcases, etc.
\usepackage{booktabs}     % \toprule etc. in tables
\usepackage{hyperref}
\usepackage{geometry}
\usepackage{microtype}
\usepackage{parskip}
\usepackage{xcolor}
\usepackage{tcolorbox}
\usepackage{enumitem}
\usepackage{float}             % [H] float placement

\geometry{a4paper, margin=2.5cm}

% -------------------------------------------------------------------------
% Theorem-like environments
% -------------------------------------------------------------------------
\theoremstyle{definition}
\newtheorem{defn}{Definition}[section]
\newtheorem{remark}{Remark}[section]

\theoremstyle{plain}
\newtheorem{prop}{Proposition}[section]

% -------------------------------------------------------------------------
% Mesh and geometry
% -------------------------------------------------------------------------
\newcommand{\V}{\mathbf{V}}            % vertex array V[i] in R^dim
\newcommand{\F}{\mathbf{F}}            % face (triangle) array
\newcommand{\vone}[1]{v_{1,#1}}       % first vertex index of VP i
\newcommand{\vtwo}[1]{v_{2,#1}}       % second vertex index of VP i

% -------------------------------------------------------------------------
% Variable points
% -------------------------------------------------------------------------
\newcommand{\lam}{\lambda}             % scalar VP parameter
\newcommand{\Lam}{\boldsymbol{\lambda}}% full parameter vector
\newcommand{\vp}[1]{\mathbf{p}_{#1}}  % 3D position of VP i
\newcommand{\ed}[1]{\mathbf{d}_{#1}}  % edge direction d_i = V[v1_i] - V[v2_i]

% -------------------------------------------------------------------------
% Steiner / triple-point
% -------------------------------------------------------------------------
\newcommand{\Sv}{\mathbf{S}}           % Steiner (Fermat-Torricelli) point
\newcommand{\Ki}{K_{i}}               % projection matrix (I-n_i n_i^T)/r_i
\newcommand{\Mmat}{M}                 % sum of K_i matrices at a triple point

% -------------------------------------------------------------------------
% Normals and unit vectors
% -------------------------------------------------------------------------
\newcommand{\nhat}{\hat{\mathbf{n}}}  % unit normal to a triangle
\newcommand{\Dhat}{\hat{\boldsymbol{\Delta}}}  % unit segment direction

% -------------------------------------------------------------------------
% Operations
% -------------------------------------------------------------------------
\newcommand{\norm}[1]{\left\|#1\right\|}
\newcommand{\abs}[1]{\left|#1\right|}
\newcommand{\ip}[2]{\langle #1,\, #2 \rangle}

% Projection perpendicular to a unit vector \uhat
\newcommand{\Proj}[1]{\bigl(\mathbf{I} - #1 #1^\top\bigr)}

% argmax operator (winner-take-all assignment in Phase 1 documents)
\DeclareMathOperator*{\argmax}{arg\,max}
\DeclareMathOperator*{\argmin}{arg\,min}

% -------------------------------------------------------------------------
% Calculus shorthands
% -------------------------------------------------------------------------
\newcommand{\pd}[2]{\dfrac{\partial #1}{\partial #2}}
\newcommand{\pdd}[3]{\dfrac{\partial^2 #1}{\partial #2\,\partial #3}}
\newcommand{\grad}{\nabla}
\newcommand{\Hess}[1]{\nabla^2 #1}

% -------------------------------------------------------------------------
% Optimization problem objects
% -------------------------------------------------------------------------
\newcommand{\Preg}{P_{\mathrm{reg}}}  % regular perimeter
\newcommand{\Pst}{P_{\mathrm{st}}}   % Steiner perimeter
\newcommand{\Areg}{A^{\mathrm{reg}}} % regular cell area
\newcommand{\Ast}{A^{\mathrm{st}}}   % Steiner area contribution
\newcommand{\Atgt}{\bar{A}}           % target area per cell
\newcommand{\Lagr}{\mathcal{L}}       % Lagrangian
\newcommand{\objfac}{\sigma}          % IPOPT objective scaling factor

% -------------------------------------------------------------------------
% Code cross-references (for "In the code..." callouts)
% -------------------------------------------------------------------------
\newcommand{\code}[1]{\texttt{#1}}
\newcommand{\codefile}[1]{\texttt{\small #1}}

% -------------------------------------------------------------------------
% Threshold dynamics / approach B (10-mbo-auction-dynamics)
% -------------------------------------------------------------------------
\newcommand{\Dlump}{D}                 % lumped mass matrix diag(v)
\newcommand{\Atau}{A_\tau}             % discrete diffusion operator (M + tau K)^{-1} M
\newcommand{\Etau}{E_\tau}             % heat-content (threshold-dynamics) energy
\newcommand{\Rcell}{R_{\mathrm{cell}}} % radius of the equal-area disc sqrt(Abar/pi)
\newcommand{\hmean}{h_{\mathrm{mean}}} % mean edge length
\newcommand{\hmax}{h_{\mathrm{max}}}   % max edge length
\newcommand{\Gt}{G_t}                  % heat kernel at time t
\newcommand{\Ktau}{K^{\mathrm{res}}_\tau} % resolvent (backward-Euler) kernel
)
% -------------------------------------------------------------------------
\usepackage{amsmath,amssymb,amsthm}
\usepackage{bm}           % \bm for bold math
\usepackage{mathtools}    % \coloneqq, dcases, etc.
\usepackage{booktabs}     % \toprule etc. in tables
\usepackage{hyperref}
\usepackage{geometry}
\usepackage{microtype}
\usepackage{parskip}
\usepackage{xcolor}
\usepackage{tcolorbox}
\usepackage{enumitem}
\usepackage{float}             % [H] float placement

\geometry{a4paper, margin=2.5cm}

% -------------------------------------------------------------------------
% Theorem-like environments
% -------------------------------------------------------------------------
\theoremstyle{definition}
\newtheorem{defn}{Definition}[section]
\newtheorem{remark}{Remark}[section]

\theoremstyle{plain}
\newtheorem{prop}{Proposition}[section]

% -------------------------------------------------------------------------
% Mesh and geometry
% -------------------------------------------------------------------------
\newcommand{\V}{\mathbf{V}}            % vertex array V[i] in R^dim
\newcommand{\F}{\mathbf{F}}            % face (triangle) array
\newcommand{\vone}[1]{v_{1,#1}}       % first vertex index of VP i
\newcommand{\vtwo}[1]{v_{2,#1}}       % second vertex index of VP i

% -------------------------------------------------------------------------
% Variable points
% -------------------------------------------------------------------------
\newcommand{\lam}{\lambda}             % scalar VP parameter
\newcommand{\Lam}{\boldsymbol{\lambda}}% full parameter vector
\newcommand{\vp}[1]{\mathbf{p}_{#1}}  % 3D position of VP i
\newcommand{\ed}[1]{\mathbf{d}_{#1}}  % edge direction d_i = V[v1_i] - V[v2_i]

% -------------------------------------------------------------------------
% Steiner / triple-point
% -------------------------------------------------------------------------
\newcommand{\Sv}{\mathbf{S}}           % Steiner (Fermat-Torricelli) point
\newcommand{\Ki}{K_{i}}               % projection matrix (I-n_i n_i^T)/r_i
\newcommand{\Mmat}{M}                 % sum of K_i matrices at a triple point

% -------------------------------------------------------------------------
% Normals and unit vectors
% -------------------------------------------------------------------------
\newcommand{\nhat}{\hat{\mathbf{n}}}  % unit normal to a triangle
\newcommand{\Dhat}{\hat{\boldsymbol{\Delta}}}  % unit segment direction

% -------------------------------------------------------------------------
% Operations
% -------------------------------------------------------------------------
\newcommand{\norm}[1]{\left\|#1\right\|}
\newcommand{\abs}[1]{\left|#1\right|}
\newcommand{\ip}[2]{\langle #1,\, #2 \rangle}

% Projection perpendicular to a unit vector \uhat
\newcommand{\Proj}[1]{\bigl(\mathbf{I} - #1 #1^\top\bigr)}

% argmax operator (winner-take-all assignment in Phase 1 documents)
\DeclareMathOperator*{\argmax}{arg\,max}
\DeclareMathOperator*{\argmin}{arg\,min}

% -------------------------------------------------------------------------
% Calculus shorthands
% -------------------------------------------------------------------------
\newcommand{\pd}[2]{\dfrac{\partial #1}{\partial #2}}
\newcommand{\pdd}[3]{\dfrac{\partial^2 #1}{\partial #2\,\partial #3}}
\newcommand{\grad}{\nabla}
\newcommand{\Hess}[1]{\nabla^2 #1}

% -------------------------------------------------------------------------
% Optimization problem objects
% -------------------------------------------------------------------------
\newcommand{\Preg}{P_{\mathrm{reg}}}  % regular perimeter
\newcommand{\Pst}{P_{\mathrm{st}}}   % Steiner perimeter
\newcommand{\Areg}{A^{\mathrm{reg}}} % regular cell area
\newcommand{\Ast}{A^{\mathrm{st}}}   % Steiner area contribution
\newcommand{\Atgt}{\bar{A}}           % target area per cell
\newcommand{\Lagr}{\mathcal{L}}       % Lagrangian
\newcommand{\objfac}{\sigma}          % IPOPT objective scaling factor

% -------------------------------------------------------------------------
% Code cross-references (for "In the code..." callouts)
% -------------------------------------------------------------------------
\newcommand{\code}[1]{\texttt{#1}}
\newcommand{\codefile}[1]{\texttt{\small #1}}

% -------------------------------------------------------------------------
% Threshold dynamics / approach B (10-mbo-auction-dynamics)
% -------------------------------------------------------------------------
\newcommand{\Dlump}{D}                 % lumped mass matrix diag(v)
\newcommand{\Atau}{A_\tau}             % discrete diffusion operator (M + tau K)^{-1} M
\newcommand{\Etau}{E_\tau}             % heat-content (threshold-dynamics) energy
\newcommand{\Rcell}{R_{\mathrm{cell}}} % radius of the equal-area disc sqrt(Abar/pi)
\newcommand{\hmean}{h_{\mathrm{mean}}} % mean edge length
\newcommand{\hmax}{h_{\mathrm{max}}}   % max edge length
\newcommand{\Gt}{G_t}                  % heat kernel at time t
\newcommand{\Ktau}{K^{\mathrm{res}}_\tau} % resolvent (backward-Euler) kernel


% -------------------------------------------------------------------------
% Packages (include here so each document only needs % shared/macros.tex
% Shared notation for all surface-partition math documents.
% Include with: % shared/macros.tex
% Shared notation for all surface-partition math documents.
% Include with: \input{../shared/macros}

% -------------------------------------------------------------------------
% Packages (include here so each document only needs \input{../shared/macros})
% -------------------------------------------------------------------------
\usepackage{amsmath,amssymb,amsthm}
\usepackage{bm}           % \bm for bold math
\usepackage{mathtools}    % \coloneqq, dcases, etc.
\usepackage{booktabs}     % \toprule etc. in tables
\usepackage{hyperref}
\usepackage{geometry}
\usepackage{microtype}
\usepackage{parskip}
\usepackage{xcolor}
\usepackage{tcolorbox}
\usepackage{enumitem}
\usepackage{float}             % [H] float placement

\geometry{a4paper, margin=2.5cm}

% -------------------------------------------------------------------------
% Theorem-like environments
% -------------------------------------------------------------------------
\theoremstyle{definition}
\newtheorem{defn}{Definition}[section]
\newtheorem{remark}{Remark}[section]

\theoremstyle{plain}
\newtheorem{prop}{Proposition}[section]

% -------------------------------------------------------------------------
% Mesh and geometry
% -------------------------------------------------------------------------
\newcommand{\V}{\mathbf{V}}            % vertex array V[i] in R^dim
\newcommand{\F}{\mathbf{F}}            % face (triangle) array
\newcommand{\vone}[1]{v_{1,#1}}       % first vertex index of VP i
\newcommand{\vtwo}[1]{v_{2,#1}}       % second vertex index of VP i

% -------------------------------------------------------------------------
% Variable points
% -------------------------------------------------------------------------
\newcommand{\lam}{\lambda}             % scalar VP parameter
\newcommand{\Lam}{\boldsymbol{\lambda}}% full parameter vector
\newcommand{\vp}[1]{\mathbf{p}_{#1}}  % 3D position of VP i
\newcommand{\ed}[1]{\mathbf{d}_{#1}}  % edge direction d_i = V[v1_i] - V[v2_i]

% -------------------------------------------------------------------------
% Steiner / triple-point
% -------------------------------------------------------------------------
\newcommand{\Sv}{\mathbf{S}}           % Steiner (Fermat-Torricelli) point
\newcommand{\Ki}{K_{i}}               % projection matrix (I-n_i n_i^T)/r_i
\newcommand{\Mmat}{M}                 % sum of K_i matrices at a triple point

% -------------------------------------------------------------------------
% Normals and unit vectors
% -------------------------------------------------------------------------
\newcommand{\nhat}{\hat{\mathbf{n}}}  % unit normal to a triangle
\newcommand{\Dhat}{\hat{\boldsymbol{\Delta}}}  % unit segment direction

% -------------------------------------------------------------------------
% Operations
% -------------------------------------------------------------------------
\newcommand{\norm}[1]{\left\|#1\right\|}
\newcommand{\abs}[1]{\left|#1\right|}
\newcommand{\ip}[2]{\langle #1,\, #2 \rangle}

% Projection perpendicular to a unit vector \uhat
\newcommand{\Proj}[1]{\bigl(\mathbf{I} - #1 #1^\top\bigr)}

% argmax operator (winner-take-all assignment in Phase 1 documents)
\DeclareMathOperator*{\argmax}{arg\,max}
\DeclareMathOperator*{\argmin}{arg\,min}

% -------------------------------------------------------------------------
% Calculus shorthands
% -------------------------------------------------------------------------
\newcommand{\pd}[2]{\dfrac{\partial #1}{\partial #2}}
\newcommand{\pdd}[3]{\dfrac{\partial^2 #1}{\partial #2\,\partial #3}}
\newcommand{\grad}{\nabla}
\newcommand{\Hess}[1]{\nabla^2 #1}

% -------------------------------------------------------------------------
% Optimization problem objects
% -------------------------------------------------------------------------
\newcommand{\Preg}{P_{\mathrm{reg}}}  % regular perimeter
\newcommand{\Pst}{P_{\mathrm{st}}}   % Steiner perimeter
\newcommand{\Areg}{A^{\mathrm{reg}}} % regular cell area
\newcommand{\Ast}{A^{\mathrm{st}}}   % Steiner area contribution
\newcommand{\Atgt}{\bar{A}}           % target area per cell
\newcommand{\Lagr}{\mathcal{L}}       % Lagrangian
\newcommand{\objfac}{\sigma}          % IPOPT objective scaling factor

% -------------------------------------------------------------------------
% Code cross-references (for "In the code..." callouts)
% -------------------------------------------------------------------------
\newcommand{\code}[1]{\texttt{#1}}
\newcommand{\codefile}[1]{\texttt{\small #1}}

% -------------------------------------------------------------------------
% Threshold dynamics / approach B (10-mbo-auction-dynamics)
% -------------------------------------------------------------------------
\newcommand{\Dlump}{D}                 % lumped mass matrix diag(v)
\newcommand{\Atau}{A_\tau}             % discrete diffusion operator (M + tau K)^{-1} M
\newcommand{\Etau}{E_\tau}             % heat-content (threshold-dynamics) energy
\newcommand{\Rcell}{R_{\mathrm{cell}}} % radius of the equal-area disc sqrt(Abar/pi)
\newcommand{\hmean}{h_{\mathrm{mean}}} % mean edge length
\newcommand{\hmax}{h_{\mathrm{max}}}   % max edge length
\newcommand{\Gt}{G_t}                  % heat kernel at time t
\newcommand{\Ktau}{K^{\mathrm{res}}_\tau} % resolvent (backward-Euler) kernel


% -------------------------------------------------------------------------
% Packages (include here so each document only needs % shared/macros.tex
% Shared notation for all surface-partition math documents.
% Include with: \input{../shared/macros}

% -------------------------------------------------------------------------
% Packages (include here so each document only needs \input{../shared/macros})
% -------------------------------------------------------------------------
\usepackage{amsmath,amssymb,amsthm}
\usepackage{bm}           % \bm for bold math
\usepackage{mathtools}    % \coloneqq, dcases, etc.
\usepackage{booktabs}     % \toprule etc. in tables
\usepackage{hyperref}
\usepackage{geometry}
\usepackage{microtype}
\usepackage{parskip}
\usepackage{xcolor}
\usepackage{tcolorbox}
\usepackage{enumitem}
\usepackage{float}             % [H] float placement

\geometry{a4paper, margin=2.5cm}

% -------------------------------------------------------------------------
% Theorem-like environments
% -------------------------------------------------------------------------
\theoremstyle{definition}
\newtheorem{defn}{Definition}[section]
\newtheorem{remark}{Remark}[section]

\theoremstyle{plain}
\newtheorem{prop}{Proposition}[section]

% -------------------------------------------------------------------------
% Mesh and geometry
% -------------------------------------------------------------------------
\newcommand{\V}{\mathbf{V}}            % vertex array V[i] in R^dim
\newcommand{\F}{\mathbf{F}}            % face (triangle) array
\newcommand{\vone}[1]{v_{1,#1}}       % first vertex index of VP i
\newcommand{\vtwo}[1]{v_{2,#1}}       % second vertex index of VP i

% -------------------------------------------------------------------------
% Variable points
% -------------------------------------------------------------------------
\newcommand{\lam}{\lambda}             % scalar VP parameter
\newcommand{\Lam}{\boldsymbol{\lambda}}% full parameter vector
\newcommand{\vp}[1]{\mathbf{p}_{#1}}  % 3D position of VP i
\newcommand{\ed}[1]{\mathbf{d}_{#1}}  % edge direction d_i = V[v1_i] - V[v2_i]

% -------------------------------------------------------------------------
% Steiner / triple-point
% -------------------------------------------------------------------------
\newcommand{\Sv}{\mathbf{S}}           % Steiner (Fermat-Torricelli) point
\newcommand{\Ki}{K_{i}}               % projection matrix (I-n_i n_i^T)/r_i
\newcommand{\Mmat}{M}                 % sum of K_i matrices at a triple point

% -------------------------------------------------------------------------
% Normals and unit vectors
% -------------------------------------------------------------------------
\newcommand{\nhat}{\hat{\mathbf{n}}}  % unit normal to a triangle
\newcommand{\Dhat}{\hat{\boldsymbol{\Delta}}}  % unit segment direction

% -------------------------------------------------------------------------
% Operations
% -------------------------------------------------------------------------
\newcommand{\norm}[1]{\left\|#1\right\|}
\newcommand{\abs}[1]{\left|#1\right|}
\newcommand{\ip}[2]{\langle #1,\, #2 \rangle}

% Projection perpendicular to a unit vector \uhat
\newcommand{\Proj}[1]{\bigl(\mathbf{I} - #1 #1^\top\bigr)}

% argmax operator (winner-take-all assignment in Phase 1 documents)
\DeclareMathOperator*{\argmax}{arg\,max}
\DeclareMathOperator*{\argmin}{arg\,min}

% -------------------------------------------------------------------------
% Calculus shorthands
% -------------------------------------------------------------------------
\newcommand{\pd}[2]{\dfrac{\partial #1}{\partial #2}}
\newcommand{\pdd}[3]{\dfrac{\partial^2 #1}{\partial #2\,\partial #3}}
\newcommand{\grad}{\nabla}
\newcommand{\Hess}[1]{\nabla^2 #1}

% -------------------------------------------------------------------------
% Optimization problem objects
% -------------------------------------------------------------------------
\newcommand{\Preg}{P_{\mathrm{reg}}}  % regular perimeter
\newcommand{\Pst}{P_{\mathrm{st}}}   % Steiner perimeter
\newcommand{\Areg}{A^{\mathrm{reg}}} % regular cell area
\newcommand{\Ast}{A^{\mathrm{st}}}   % Steiner area contribution
\newcommand{\Atgt}{\bar{A}}           % target area per cell
\newcommand{\Lagr}{\mathcal{L}}       % Lagrangian
\newcommand{\objfac}{\sigma}          % IPOPT objective scaling factor

% -------------------------------------------------------------------------
% Code cross-references (for "In the code..." callouts)
% -------------------------------------------------------------------------
\newcommand{\code}[1]{\texttt{#1}}
\newcommand{\codefile}[1]{\texttt{\small #1}}

% -------------------------------------------------------------------------
% Threshold dynamics / approach B (10-mbo-auction-dynamics)
% -------------------------------------------------------------------------
\newcommand{\Dlump}{D}                 % lumped mass matrix diag(v)
\newcommand{\Atau}{A_\tau}             % discrete diffusion operator (M + tau K)^{-1} M
\newcommand{\Etau}{E_\tau}             % heat-content (threshold-dynamics) energy
\newcommand{\Rcell}{R_{\mathrm{cell}}} % radius of the equal-area disc sqrt(Abar/pi)
\newcommand{\hmean}{h_{\mathrm{mean}}} % mean edge length
\newcommand{\hmax}{h_{\mathrm{max}}}   % max edge length
\newcommand{\Gt}{G_t}                  % heat kernel at time t
\newcommand{\Ktau}{K^{\mathrm{res}}_\tau} % resolvent (backward-Euler) kernel
)
% -------------------------------------------------------------------------
\usepackage{amsmath,amssymb,amsthm}
\usepackage{bm}           % \bm for bold math
\usepackage{mathtools}    % \coloneqq, dcases, etc.
\usepackage{booktabs}     % \toprule etc. in tables
\usepackage{hyperref}
\usepackage{geometry}
\usepackage{microtype}
\usepackage{parskip}
\usepackage{xcolor}
\usepackage{tcolorbox}
\usepackage{enumitem}
\usepackage{float}             % [H] float placement

\geometry{a4paper, margin=2.5cm}

% -------------------------------------------------------------------------
% Theorem-like environments
% -------------------------------------------------------------------------
\theoremstyle{definition}
\newtheorem{defn}{Definition}[section]
\newtheorem{remark}{Remark}[section]

\theoremstyle{plain}
\newtheorem{prop}{Proposition}[section]

% -------------------------------------------------------------------------
% Mesh and geometry
% -------------------------------------------------------------------------
\newcommand{\V}{\mathbf{V}}            % vertex array V[i] in R^dim
\newcommand{\F}{\mathbf{F}}            % face (triangle) array
\newcommand{\vone}[1]{v_{1,#1}}       % first vertex index of VP i
\newcommand{\vtwo}[1]{v_{2,#1}}       % second vertex index of VP i

% -------------------------------------------------------------------------
% Variable points
% -------------------------------------------------------------------------
\newcommand{\lam}{\lambda}             % scalar VP parameter
\newcommand{\Lam}{\boldsymbol{\lambda}}% full parameter vector
\newcommand{\vp}[1]{\mathbf{p}_{#1}}  % 3D position of VP i
\newcommand{\ed}[1]{\mathbf{d}_{#1}}  % edge direction d_i = V[v1_i] - V[v2_i]

% -------------------------------------------------------------------------
% Steiner / triple-point
% -------------------------------------------------------------------------
\newcommand{\Sv}{\mathbf{S}}           % Steiner (Fermat-Torricelli) point
\newcommand{\Ki}{K_{i}}               % projection matrix (I-n_i n_i^T)/r_i
\newcommand{\Mmat}{M}                 % sum of K_i matrices at a triple point

% -------------------------------------------------------------------------
% Normals and unit vectors
% -------------------------------------------------------------------------
\newcommand{\nhat}{\hat{\mathbf{n}}}  % unit normal to a triangle
\newcommand{\Dhat}{\hat{\boldsymbol{\Delta}}}  % unit segment direction

% -------------------------------------------------------------------------
% Operations
% -------------------------------------------------------------------------
\newcommand{\norm}[1]{\left\|#1\right\|}
\newcommand{\abs}[1]{\left|#1\right|}
\newcommand{\ip}[2]{\langle #1,\, #2 \rangle}

% Projection perpendicular to a unit vector \uhat
\newcommand{\Proj}[1]{\bigl(\mathbf{I} - #1 #1^\top\bigr)}

% argmax operator (winner-take-all assignment in Phase 1 documents)
\DeclareMathOperator*{\argmax}{arg\,max}
\DeclareMathOperator*{\argmin}{arg\,min}

% -------------------------------------------------------------------------
% Calculus shorthands
% -------------------------------------------------------------------------
\newcommand{\pd}[2]{\dfrac{\partial #1}{\partial #2}}
\newcommand{\pdd}[3]{\dfrac{\partial^2 #1}{\partial #2\,\partial #3}}
\newcommand{\grad}{\nabla}
\newcommand{\Hess}[1]{\nabla^2 #1}

% -------------------------------------------------------------------------
% Optimization problem objects
% -------------------------------------------------------------------------
\newcommand{\Preg}{P_{\mathrm{reg}}}  % regular perimeter
\newcommand{\Pst}{P_{\mathrm{st}}}   % Steiner perimeter
\newcommand{\Areg}{A^{\mathrm{reg}}} % regular cell area
\newcommand{\Ast}{A^{\mathrm{st}}}   % Steiner area contribution
\newcommand{\Atgt}{\bar{A}}           % target area per cell
\newcommand{\Lagr}{\mathcal{L}}       % Lagrangian
\newcommand{\objfac}{\sigma}          % IPOPT objective scaling factor

% -------------------------------------------------------------------------
% Code cross-references (for "In the code..." callouts)
% -------------------------------------------------------------------------
\newcommand{\code}[1]{\texttt{#1}}
\newcommand{\codefile}[1]{\texttt{\small #1}}

% -------------------------------------------------------------------------
% Threshold dynamics / approach B (10-mbo-auction-dynamics)
% -------------------------------------------------------------------------
\newcommand{\Dlump}{D}                 % lumped mass matrix diag(v)
\newcommand{\Atau}{A_\tau}             % discrete diffusion operator (M + tau K)^{-1} M
\newcommand{\Etau}{E_\tau}             % heat-content (threshold-dynamics) energy
\newcommand{\Rcell}{R_{\mathrm{cell}}} % radius of the equal-area disc sqrt(Abar/pi)
\newcommand{\hmean}{h_{\mathrm{mean}}} % mean edge length
\newcommand{\hmax}{h_{\mathrm{max}}}   % max edge length
\newcommand{\Gt}{G_t}                  % heat kernel at time t
\newcommand{\Ktau}{K^{\mathrm{res}}_\tau} % resolvent (backward-Euler) kernel
)
% -------------------------------------------------------------------------
\usepackage{amsmath,amssymb,amsthm}
\usepackage{bm}           % \bm for bold math
\usepackage{mathtools}    % \coloneqq, dcases, etc.
\usepackage{booktabs}     % \toprule etc. in tables
\usepackage{hyperref}
\usepackage{geometry}
\usepackage{microtype}
\usepackage{parskip}
\usepackage{xcolor}
\usepackage{tcolorbox}
\usepackage{enumitem}
\usepackage{float}             % [H] float placement

\geometry{a4paper, margin=2.5cm}

% -------------------------------------------------------------------------
% Theorem-like environments
% -------------------------------------------------------------------------
\theoremstyle{definition}
\newtheorem{defn}{Definition}[section]
\newtheorem{remark}{Remark}[section]

\theoremstyle{plain}
\newtheorem{prop}{Proposition}[section]

% -------------------------------------------------------------------------
% Mesh and geometry
% -------------------------------------------------------------------------
\newcommand{\V}{\mathbf{V}}            % vertex array V[i] in R^dim
\newcommand{\F}{\mathbf{F}}            % face (triangle) array
\newcommand{\vone}[1]{v_{1,#1}}       % first vertex index of VP i
\newcommand{\vtwo}[1]{v_{2,#1}}       % second vertex index of VP i

% -------------------------------------------------------------------------
% Variable points
% -------------------------------------------------------------------------
\newcommand{\lam}{\lambda}             % scalar VP parameter
\newcommand{\Lam}{\boldsymbol{\lambda}}% full parameter vector
\newcommand{\vp}[1]{\mathbf{p}_{#1}}  % 3D position of VP i
\newcommand{\ed}[1]{\mathbf{d}_{#1}}  % edge direction d_i = V[v1_i] - V[v2_i]

% -------------------------------------------------------------------------
% Steiner / triple-point
% -------------------------------------------------------------------------
\newcommand{\Sv}{\mathbf{S}}           % Steiner (Fermat-Torricelli) point
\newcommand{\Ki}{K_{i}}               % projection matrix (I-n_i n_i^T)/r_i
\newcommand{\Mmat}{M}                 % sum of K_i matrices at a triple point

% -------------------------------------------------------------------------
% Normals and unit vectors
% -------------------------------------------------------------------------
\newcommand{\nhat}{\hat{\mathbf{n}}}  % unit normal to a triangle
\newcommand{\Dhat}{\hat{\boldsymbol{\Delta}}}  % unit segment direction

% -------------------------------------------------------------------------
% Operations
% -------------------------------------------------------------------------
\newcommand{\norm}[1]{\left\|#1\right\|}
\newcommand{\abs}[1]{\left|#1\right|}
\newcommand{\ip}[2]{\langle #1,\, #2 \rangle}

% Projection perpendicular to a unit vector \uhat
\newcommand{\Proj}[1]{\bigl(\mathbf{I} - #1 #1^\top\bigr)}

% argmax operator (winner-take-all assignment in Phase 1 documents)
\DeclareMathOperator*{\argmax}{arg\,max}
\DeclareMathOperator*{\argmin}{arg\,min}

% -------------------------------------------------------------------------
% Calculus shorthands
% -------------------------------------------------------------------------
\newcommand{\pd}[2]{\dfrac{\partial #1}{\partial #2}}
\newcommand{\pdd}[3]{\dfrac{\partial^2 #1}{\partial #2\,\partial #3}}
\newcommand{\grad}{\nabla}
\newcommand{\Hess}[1]{\nabla^2 #1}

% -------------------------------------------------------------------------
% Optimization problem objects
% -------------------------------------------------------------------------
\newcommand{\Preg}{P_{\mathrm{reg}}}  % regular perimeter
\newcommand{\Pst}{P_{\mathrm{st}}}   % Steiner perimeter
\newcommand{\Areg}{A^{\mathrm{reg}}} % regular cell area
\newcommand{\Ast}{A^{\mathrm{st}}}   % Steiner area contribution
\newcommand{\Atgt}{\bar{A}}           % target area per cell
\newcommand{\Lagr}{\mathcal{L}}       % Lagrangian
\newcommand{\objfac}{\sigma}          % IPOPT objective scaling factor

% -------------------------------------------------------------------------
% Code cross-references (for "In the code..." callouts)
% -------------------------------------------------------------------------
\newcommand{\code}[1]{\texttt{#1}}
\newcommand{\codefile}[1]{\texttt{\small #1}}

% -------------------------------------------------------------------------
% Threshold dynamics / approach B (10-mbo-auction-dynamics)
% -------------------------------------------------------------------------
\newcommand{\Dlump}{D}                 % lumped mass matrix diag(v)
\newcommand{\Atau}{A_\tau}             % discrete diffusion operator (M + tau K)^{-1} M
\newcommand{\Etau}{E_\tau}             % heat-content (threshold-dynamics) energy
\newcommand{\Rcell}{R_{\mathrm{cell}}} % radius of the equal-area disc sqrt(Abar/pi)
\newcommand{\hmean}{h_{\mathrm{mean}}} % mean edge length
\newcommand{\hmax}{h_{\mathrm{max}}}   % max edge length
\newcommand{\Gt}{G_t}                  % heat kernel at time t
\newcommand{\Ktau}{K^{\mathrm{res}}_\tau} % resolvent (backward-Euler) kernel


\hypersetup{
  colorlinks = true,
  linkcolor  = blue!60!black,
  citecolor  = green!50!black,
  urlcolor   = blue!60!black,
  pdftitle   = {Auction-Dynamics MBO on Surface Meshes},
  pdfauthor  = {surface-partition project}
}

\title{\textbf{Auction-Dynamics MBO on Surface Meshes}\\[0.5em]
  \large The threshold-dynamics replacement for Phase~1: provenance, derivation,
  what transfers from the literature, and what does not}
\author{}
\date{September 2026}

\begin{document}
\maketitle

\begin{abstract}
Approach~B (\codefile{src/partition/mbo\_auction.py}) replaces the Phase~1
$\Gamma$-convergence relaxation by volume-constrained multiphase threshold
dynamics: the state is a hard labelling of the mesh vertices, each step diffuses
every cell's indicator by one backward-Euler step of the surface heat equation and
reassigns every vertex by a balanced thresholding $\argmax_k [y_{ik} + \psi_k]$
whose offsets $\psi$ enforce equal discrete area. None of the ingredients is new:
the scheme is Merriman--Bence--Osher's, its variational interpretation is
Esedo\u{g}lu--Otto's, the balanced thresholding is Jacobs--Merkurjev--Esedo\u{g}lu's
auction dynamics, and the implicit-Euler diffusion is the one the graph-MBO
literature uses. This document first records, from the papers themselves, what
each of those works did, at what cell counts, on what domains --- because an
earlier description of this approach overstated its novelty on three of its
four clauses. It then derives what is specific to this implementation: that the
backward-Euler step is an \emph{exponential mixture} of heat kernels, from which
its interface constant ($\tfrac12$, against $1/\sqrt\pi$ for the Gaussian) and its
per-step displacement ($\tau\kappa/2$, against $\tau\kappa$) follow in three
lines each and are confirmed on the production mesh; that the balanced
threshold is the dual of a transportation problem whose two-phase case is the
Ruuth--Wetton shifted level; that the Esedo\u{g}lu--Otto dissipation theorem does
\emph{not} apply to the implemented pairing of a consistent diffusion with a
lumped assignment and an inexact dual, and exactly which inequality survives;
that a lumped--lumped pairing would restore the theorem (not implemented); and
the two-sided time-step window, with the over-merge side quantified by the
diffused indicator's value at a cell's centre, $1 - x K_1(x)$ at
$x = \Rcell/\sqrt\tau$. Measurements are cited from report~08; nothing here is a
new measurement except the one-step displacement check.
\end{abstract}

\tableofcontents
\newpage

\section{Overview}
\label{sec:overview}

\begin{tcolorbox}[colback=gray!6!white, colframe=gray!50!black,
                  title={Provenance of every computed number in this document},
                  fontupper=\small]
\begin{tabular}{@{}p{2.6cm}p{11.8cm}@{}}
Status & derivation, complete; twice adversarially reviewed (2026-09-17), all
  findings applied \\
Script & \codefile{docs/math/10-mbo-auction-dynamics/check\_numerics.py}
  $\to$ committed \codefile{numerics.yaml} beside this file. \textbf{Every number
  below that was produced by running code is quoted from that file}, with its
  key named in the text; nothing is quoted from an interactive session or from
  a reviewer's scratch script. Run from the repo root (\(\approx 5\) min;
  \code{-{}-quick} skips the two \(100 \times 96\) dense eigendecompositions). \\
Determinism & jittered disc centres (\S\ref{sec:check}):
  \code{numpy.random.default\_rng(0)}, 24 samples per row, one shared generator
  consumed row by row in the table's order, offsets
  \(\theta \sim U(0, 2\pi/348)\), \(\varphi \sim U(0, 2\pi/328)\). Labelling pair
  (\S\ref{sec:outside}): the geodesic balanced init at the project seed
  \(84172851\), then one full-budget balanced step. \\
Meshes & torus \(R = 1.0\), \(r = 0.6\) throughout (\code{TorusMeshProvider});
  \(40 \times 36\), \(60 \times 52\), \(100 \times 96\) (the production base), and
  the \(348 \times 328\) deliverable mesh (\(V = 114{,}144\)); \(\tau\) from
  \code{tau\_diagnostics} with \code{MBOConfig()} defaults. \\
Versions & Python 3.12.13, numpy 2.4.4, scipy 1.17.1; macOS 15.3.1 arm64
  (Mac mini). Eigenvalues may differ in the last printed digit across BLAS
  builds; nothing quoted depends on that digit. \\
Cited, not computed & the run statistics of \S\ref{sec:outside} and
  \S\ref{sec:survives} (4 of 96 energy increases; the slack range) are from
  \codefile{docs/experiments/08-mbo-auction-dynamics/data.yaml}; the
  \(68.16\%\)/\(1.97\%\) of \S\ref{sec:init} from
  \codefile{docs/plans/PHASE1\_BC\_REPLACEMENT\_PLAN.md}; the freeze constants'
  \(c = 2\)/\(c = 4\) pair from \codefile{testing/test\_mbo\_auction.py}; every
  statement about a paper from the copy indexed in
  \codefile{docs/papers/README.md}. \\
\end{tabular}
\end{tcolorbox}

Phase~1 of this project, as specified by Bogosel and Oudet
\cite{bogosel2017partitions} and derived in document~06
(\codefile{docs/math/06-phase1-energy-discretization/}), minimises a
Modica--Mortola functional over $N$ continuous densities under an equal-mass
constraint, and reads the partition off as the winner-take-all territory of those
densities. Document~09 (\codefile{docs/math/09-balanced-readout/}) derives why
that readout can diverge from what was constrained --- at high $N$ it
manufactures cells whose territory is far from the equal-area target, and cells
whose territory is disconnected --- and the extraction-time correction that
repairs it.

Approach~B removes the category of problem rather than repairing it. Its state is
a labelling $\omega : \{1,\dots,V\} \to \{1,\dots,N\}$ of the mesh vertices, never
a continuous field, so there is no readout to diverge. Each step at a fixed mesh
level is
\begin{enumerate}[itemsep=1pt]
  \item \textbf{diffuse}: for every cell $k$, solve $(M + \tau K)\, y_k = M \chi_k$,
    where $\chi_k$ is the indicator of cell $k$'s territory;
  \item \textbf{balanced threshold}: $\omega'(i) = \argmax_k [\, y_{ik} + \psi_k \,]$,
    with the $N$ offsets $\psi$ chosen so that every cell's discrete area equals
    the target $\Atgt$.
\end{enumerate}
This is repeated until the labels stop changing, the labels are carried to the
next finer mesh by nearest-vertex interpolation, and the process repeats. Phase~2
then refines the resulting contour exactly as it would a Phase~1 output
(\codefile{src/optimization/perimeter\_optimizer.py}); nothing downstream
changed.

The measurements --- perimeter against the validated PGD deliverables, the three
validity gates, attribution to the dynamics rather than the initialisation, seed
spread, and the $N = 400$--$1000$ runs --- are in
\codefile{docs/experiments/08-mbo-auction-dynamics/} and are not repeated here.
This document is the derivation and the provenance. Following the scope policy of
\codefile{docs/math/AUTHORING\_GUIDE.md}, it derives only what
\codefile{src/partition/mbo\_auction.py} computes; where a cleaner variant exists
on paper (\S\ref{sec:lumped-pairing}) it is stated and marked as not implemented.

\section{Provenance: what was done before, and what is claimed here}
\label{sec:provenance}

Every reference in this section was read from the copy held in
\codefile{docs/papers/} (title page, abstract, and the sections cited), not from a
secondary description.

\subsection{The lineage}
\label{sec:lineage}

\paragraph{Threshold dynamics.} Merriman, Bence and Osher introduced the scheme in
a 1992 UCLA report \cite{merriman1992diffusion}: diffuse the indicator of a set
for a short time, threshold at $\tfrac12$, repeat; the boundary moves by mean
curvature. Their 1994 paper \cite{merriman1994motion} extends it to several phases
by diffusing every indicator and assigning each point to the largest --- the
$\argmax$ step above. Convergence of the two-phase scheme to mean-curvature flow
was proved by Evans \cite{evans1993convergence}; the multiphase case, in the
sense of a BV weak formulation and conditionally on energy convergence, by Laux
and Otto \cite{laux2016convergence}. Ruuth \cite{ruuth1998multiphase} extended
the multiphase scheme to prescribed junction angles.

\paragraph{The variational interpretation.} Esedo\u{g}lu and Otto
\cite{esedoglu2015threshold} showed that one MBO step is the minimiser of the
\emph{linearisation} of a ``heat content'' energy $\Etau$ at the current
configuration, that this energy is concave on the convex hull of partitions, so
the step never increases it (their Lemma~5.2, for any kernel $G \ge 0$ with
$\widehat G \ge 0$, \S5.2), and that $\Etau$ $\Gamma$-converges to the total
interfacial energy as $\tau \to 0$ (their Appendix~A, for smooth radial kernels).
Esedo\u{g}lu and Jacobs \cite{esedoglu2017kernels} widen the $\Gamma$-convergence
to kernels that are merely nonnegative, symmetric, integrable with a finite first
moment, and positive near the origin (their Theorem~5.1), which is the class the
present kernel belongs to (\S\ref{sec:kernel}). The underlying nonlocal-to-local $\Gamma$-limit
theory --- the convergence of a nonlocal perimeter approximation to the
surface-tension model, which \cite[ref.~1]{esedoglu2015threshold} invokes for
``heat content tends to perimeter'' --- is Alberti and Bellettini
\cite{alberti1998nonlocal}.

\paragraph{Volume constraints.} Ruuth and Wetton \cite{ruuth2003volume} enforce a
prescribed volume in the two-phase scheme by thresholding the diffused
indicator not at $\tfrac12$ but at the level $l$ that preserves the phase volume
(their Algorithm, step 2(b)); \cite{laux2017bulk} proves the scheme's
convergence. Jacobs, Merkurjev and Esedo\u{g}lu
\cite{jacobs2018auction} generalise this to $N$ phases: minimising the
linearisation of $\Etau$ under volume constraints is a linear assignment problem,
its Lagrange multipliers shift the thresholds (their eq.~(11), which is exactly
the $\argmax_k[y_{ik} + \psi_k]$ of this implementation), and they solve it with
Bertsekas' auction algorithm \cite{bertsekas1988auction}; hence ``auction
dynamics''. Laux and Swartz \cite{laux2017bulk} prove convergence of the
volume-preserving thresholding scheme.

\paragraph{Graphs, and the implicit-Euler diffusion.} Garcia-Cardona et al.\
\cite{garcia2014multiclass} run multiclass MBO on a weighted graph and replace
the Gaussian convolution by an \emph{implicit-Euler} step of the graph heat
equation, $U^{n+1/2} = B^{-1}[\,U^n - \delta t\, \mu (U^n - \hat U)\,]$ with
$B = I + \delta t\, L_s$ (their eqs.~(27)--(28); eq.~(29) is $B$ in a truncated
eigenbasis). Jacobs et al.\ \cite[\S5]{jacobs2018auction} run
\emph{volume-constrained} auction dynamics on an unstructured $k$-nearest-neighbour
graph with unit-volume nodes --- but with the symmetrised similarity matrix
$W^\top W$ used \emph{directly} as the kernel, no Laplacian and no implicit
step. Van Gennip, Guillen, Osting and Bertozzi
\cite{vangennip2014mean} prove the graph scheme is \emph{pinned} for $\tau$ below a
spectral bound and \emph{trivial} above another (their Theorems~4.2--4.4) --- the
rigorous form of the two-sided window of \S\ref{sec:tau}.

\paragraph{Surfaces.} Merriman and Ruuth \cite{merriman2007surfaces} move curves
with junctions on a curved surface by extending each indicator off the surface
with the closest-point map, diffusing in $\mathbb{R}^3$ on a Cartesian grid, and
thresholding on the surface; their \S8.3 evolves five regions on an embedded
torus and, in their Fig.~11, the 174 countries plus ocean of a world map on the
sphere --- many-region multiphase threshold dynamics on a curved closed surface,
unconstrained. Wang and Osting \cite{wang2019dirichlet} compute \emph{Dirichlet}
partitions (minimal sum of first Laplace--Beltrami eigenvalues, not perimeter) of
flat tori and of the sphere by a diffusion-generated method, via FFT and the
spherical harmonic transform, for $k \le 20$. The surface finite-element setting
used here --- P1 mass and stiffness matrices on a triangulation of an embedded
surface --- is Dziuk's \cite{dziuk1988beltrami, dziuk2013surfacefem}.

\paragraph{Partitions with many cells.} Elsey, Esedo\u{g}lu and Smereka simulate
grain growth by diffusion-generated motion with $166{,}927$ grains in two
dimensions \cite{elsey2009grain} and $133{,}110$ in three \cite{elsey2011large},
on uniform grids, without volume constraints. Wang et al.\
\cite{wang2017thresholding} apply iterative convolution--thresholding to
piecewise-constant Mumford--Shah segmentation with a handful of phases; Hu, Liu
and Wang \cite{hu2024iterative} use auction dynamics as the inner solver of the
\emph{outer} (P\'olya) problem on $256^2$ and $128^3$ grids with $n \le 9$ cells.
Bogosel and Oudet \cite{bogosel2021longest} initialise the same outer problem
from capacity-constrained Voronoi diagrams with more than $100$ equal-area cells
in the plane.

\begin{table}[H]
\centering\footnotesize
\begin{tabular}{@{}p{2.5cm}p{3.5cm}p{1.9cm}p{1.9cm}p{4.4cm}@{}}
\toprule
Work & Domain, discretisation & Cells & Equal-area & Relation to this implementation \\
\midrule
\cite{jacobs2018auction} \S4.2 & flat 2-torus $[0,1]^2$, grid, FFT & 17, \textbf{64}; 160 & \textbf{yes} (160: preserved) & same step, same problem class, exact auction \\
\cite{jacobs2018auction} \S4.2 & flat 3-torus & 8; 32 & yes (32: preserved) & Weaire--Phelan found with temperature \\
\cite{merriman2007surfaces} \S8.3 & embedded torus and sphere, closest point, grid & 5 (torus); 175 (sphere) & no & multiphase MBO on curved closed surfaces \\
\cite{wang2019dirichlet} & flat tori and $S^2$, FFT / SHT & $\le 20$ & no & other energy; curved closed surface \\
\cite{hu2024iterative} & box, $256^2$ / $128^3$ grid & $\le 9$ & yes & auction dynamics as inner solver \\
\cite{garcia2014multiclass} & weighted graph, eigenbasis & $\sim 10$ classes & no & implicit-Euler diffusion step \\
\cite{jacobs2018auction} \S5 & unstructured $k$NN graph, kernel $W^\top W$ & $\sim 10$ classes & bounds & constrained + unstructured, no surface \\
\cite{elsey2009grain, elsey2011large} & grids & $10^5$ & no & scale, unconstrained \\
\cite{bogosel2021longest} & plane & $> 100$ (init only) & yes & capacity-constrained Voronoi as init \\
\cite{bogosel2017partitions} & meshed surfaces, $\Gamma$-convergence & $\le 11$ torus, $\le 32$ sphere & yes & the Phase~1 incumbent; no thresholding \\
\bottomrule
\end{tabular}
\caption{Cell counts and settings, read from the cited sections.
``Preserved'' means each Voronoi cell keeps its \emph{initial} area or volume
(160 cells in 2D, 32 grains in 3D), not that they are equal; ``init only''
means the count appears in an initialisation figure, not a minimisation run.}
\label{tab:lineage}
\end{table}

\subsection{The corrected statement of what is new}
\label{sec:novelty}

An earlier description of approach~B (the module docstring, and \S2 of
report~08 as first written) said that every published implementation of this
machinery runs ``on a uniform Cartesian grid via FFT, on a flat domain, at small
$N$, for P\'olya's outer problem''. Table~\ref{tab:lineage} falsifies three of
those four clauses: \cite{jacobs2018auction} \S4.2 solves the \emph{inner}
problem at $N = 64$ and $160$ (two clauses), and \cite{merriman2007surfaces}
runs the multiphase scheme on a curved torus and on the sphere (the ``flat''
clause; JME's torus is flat). Only ``uniform grid via FFT'' survives, and only
for the Euclidean implementations; the graph line never used an FFT.

What the held literature does \emph{not} contain, stated as ``we are not aware
of'' and not as ``first'':
\begin{enumerate}[itemsep=1pt]
  \item the scheme on a \textbf{curved, embedded surface discretised by an
    unstructured triangle mesh}, with the diffusion done by a surface-FEM
    backward-Euler step --- the graph line has the operator form, the closest-point
    line has the surface, neither has both, and neither works with lumped and
    consistent mass on a triangulation;
  \item \textbf{equal-area constraints at $N = 400$--$1000$}: in the held
    papers, equal-area minimisation reaches 64 cells
    \cite[\S4.2]{jacobs2018auction}; the 160 and 32 there are area-preserving,
    and the ``more than 100'' of \cite{bogosel2021longest} is an initialisation
    figure;
  \item the assignment solved by an \textbf{inexact} transportation dual
    (\S\ref{sec:assignment}) rather than an exact auction, together with the
    inequality that survives inexactness (\S\ref{sec:survives});
  \item the \textbf{multilevel mesh ladder} with label interpolation and a
    per-level, mesh-derived time step (\S\ref{sec:ladder}, \S\ref{sec:tau});
  \item coupling to an exact contour refinement (Phase~2) and the gate
    instruments of report~07.
\end{enumerate}
The defensible framing is therefore \emph{combination, surface, and scale}, not
\emph{method}. The method is \cite{jacobs2018auction}'s; the diffusion operator
is the graph-MBO one transplanted to surface finite elements; ``large $N$'' by
itself is not new in three separate lines of work. And the held papers come
closer than a one-line summary suggests: \cite[\S5]{jacobs2018auction} already
has the volume constraint on an \emph{unstructured} domain (a similarity graph),
\cite[Fig.~11]{merriman2007surfaces} has 175 regions on a curved closed surface,
and \cite{wang2019dirichlet} has a curved closed surface with exact spectral
diffusion at 20 cells. What none of them has is all of it at once on a
triangulated surface with equal-area constraints in the hundreds.

\section{Notation and Setup}
\label{sec:notation}

The surface $S$ is closed and triangulated with vertices $\V = \{x_1, \dots, x_V\}$
and faces $\F$. $M$ and $K$ are the P1 mass and stiffness matrices of
\codefile{docs/math/06-phase1-energy-discretization/} eq.~(MK-def), assembled by
\code{TriMesh.compute\_matrices}; $K$ is the cotangent Laplacian and $K\mathbf 1 = 0$.
The lumped mass is $v_i = \sum_j M_{ij} > 0$, $\Dlump = \operatorname{diag}(v)$,
$|S| = \sum_i v_i$, and the equal-area target is $\Atgt = |S|/N$. The mesh scales
are the mean and maximum edge length $\hmean$, $\hmax$ over all triangle sides,
and the cell scale is the radius of the equal-area disc,
\begin{equation}
  \Rcell \;=\; \sqrt{\Atgt/\pi}.
  \label{eq:rcell}
\end{equation}

A labelling $\omega$ has indicator columns $\chi_k \in \{0,1\}^V$,
$\chi_{ik} = [\omega(i) = k]$, so $\sum_k \chi_k = \mathbf 1$, and territories
$\Omega_k = \{i : \omega(i) = k\}$ with discrete areas
$T_k = \sum_{i \in \Omega_k} v_i = v^\top \chi_k$. The diffused indicators are the
columns $y_k$ of $y \in \mathbb{R}^{V \times N}$. In the continuum, $\Gt$ is the
heat kernel of $S$ at time $t$ (on $\mathbb{R}^2$,
$\Gt(x) = (4\pi t)^{-1} e^{-|x|^2/4t}$), and $\Ktau$ the kernel of
$(I - \tau\Delta)^{-1}$, defined in \S\ref{sec:kernel}.

\begin{tcolorbox}[colback=blue!4!white, colframe=blue!40!black,
                  title={Code: the objects}, fontupper=\small]
\codefile{src/partition/mbo\_auction.py}: \code{sparse\_one\_hot(labels, n\_cells)}
builds $\chi$; \code{edge\_lengths(mesh)} gives $\hmean$, $\hmax$;
\code{cell\_radius(total\_area, n\_cells)} is \eqref{eq:rcell}; \code{mesh.v} is
$v$, \code{mesh.M}, \code{mesh.K} are $M$, $K$.
\end{tcolorbox}

\section{Threshold dynamics in the continuum}
\label{sec:continuum}

\subsection{The MBO scheme}

For a set $\Sigma \subset S$ and a time step $\tau$, one step of
\cite{merriman1992diffusion} is
\begin{equation}
  \Sigma' \;=\; \bigl\{ x : (G_\tau * \chi_\Sigma)(x) \ge \tfrac12 \bigr\},
  \label{eq:mbo2}
\end{equation}
and for a partition $(\Omega_1, \dots, \Omega_N)$ one step of
\cite{merriman1994motion} is
\begin{equation}
  \omega'(x) \;=\; \argmax_k \, (G_\tau * \chi_{\Omega_k})(x).
  \label{eq:mboN}
\end{equation}
Since $\sum_k G_\tau * \chi_{\Omega_k} = 1$, \eqref{eq:mboN} reduces to
\eqref{eq:mbo2} when $N = 2$. Away from junctions, each interface moves by mean
curvature over time $\tau$ per step, to leading order in $\tau$
\cite{evans1993convergence, esedoglu2015threshold}; \S\ref{sec:displacement}
derives the constant for the kernel actually used here.

\subsection{The heat-content energy and the linearisation principle}
\label{sec:heat-content}

Following \cite[\S4]{esedoglu2015threshold}, define for a partition
\begin{equation}
  \Etau(\omega) \;=\; \frac{1}{\sqrt\tau} \sum_{k=1}^{N}
    \int_S \chi_{\Omega_k}\,\bigl(1 - G_\tau * \chi_{\Omega_k}\bigr)\, dA
  \;=\; \frac{1}{\sqrt\tau} \sum_{k \ne j} \int_S \chi_{\Omega_k}\, G_\tau * \chi_{\Omega_j}\, dA,
  \label{eq:heat-content}
\end{equation}
the second form by $\sum_j G_\tau * \chi_{\Omega_j} = 1$. Each term measures the
heat that escapes from cell $k$ into the others in time $\tau$ --- ``heat
content''. Note the sum is over \emph{ordered} pairs, so every interface is
counted from both sides; this is eq.~(4.2) of \cite{esedoglu2015threshold} with
all surface tensions equal to one and the same $\tau^{-1/2}$ normalisation.

Relaxed to $u : S \to \Delta_N$ (the probability simplex), the energy is
\begin{equation}
  \Etau(u) \;=\; \frac{1}{\sqrt\tau}\Bigl[\, |S| \;-\; \sum_k \int_S u_k\, G_\tau * u_k \, dA \Bigr],
  \label{eq:relaxed}
\end{equation}
which is \textbf{concave}: $u_k \mapsto \int u_k\, G_\tau * u_k$ is a quadratic
form with a positive semi-definite kernel ($\widehat{G_\tau} = e^{-\tau|\xi|^2} \ge 0$),
so its negative is concave. Its linearisation at $u$ is, up to a constant,
\begin{equation}
  L_u(u') \;=\; -\frac{2}{\sqrt\tau} \sum_k \int_S u'_k\, G_\tau * u_k\, dA,
  \label{eq:linearisation}
\end{equation}
and minimising $L_u$ over simplex-valued $u'$ is a pointwise problem whose
solution puts all of $u'(x)$ on the largest $(G_\tau * u_k)(x)$ --- i.e.\ exactly
the step \eqref{eq:mboN}. This is the observation of
\cite{esedoglu2015threshold} that makes threshold dynamics a
minimising-movements scheme, and it carries the dissipation property:

\begin{prop}[Dissipation; {\cite[Lemma~5.2]{esedoglu2015threshold}}]
\label{prop:dissipation}
Let $F$ be concave and differentiable on a closed convex set $\mathcal K$, and
let $u' \in \argmin_{\mathcal K} L_u$ where $L_u$ is the linearisation of $F$ at
$u \in \mathcal K$. Then $F(u') \le F(u)$.
\end{prop}
\begin{proof}
Concavity gives $F(u') \le L_u(u')$; optimality of $u'$ and feasibility of $u$
give $L_u(u') \le L_u(u) = F(u)$.
\end{proof}

The proof uses only that $u$ itself lies in $\mathcal K$. So the conclusion
survives \emph{any} further convex constraint on $u'$ that $u$ already satisfies
--- in particular the equal-area constraint of \S\ref{sec:assignment}, provided
the constrained linearisation is minimised \emph{exactly}. That is the route by
which \cite{jacobs2018auction} obtain dissipation for auction dynamics, and
\S\ref{sec:monotonicity} examines which of its hypotheses the implementation
violates.

\subsection{\texorpdfstring{$\Gamma$}{Gamma}-convergence and the interface constant}
\label{sec:gamma}

For a kernel $K_\tau(x) = \tau^{-d/2} K(x/\sqrt\tau)$ with $K \ge 0$, radial,
$\int K = 1$ and a finite first moment, the heat content of a smooth set
satisfies, per unit interface length and to leading order in $\tau$,
\begin{equation}
  \frac{1}{\sqrt\tau} \int_{\Sigma^c} K_\tau * \chi_\Sigma \, dA
  \;\longrightarrow\; c_K \, \operatorname{Per}(\Sigma),
  \qquad
  c_K \;=\; \frac12 \int_{\mathbb{R}^d} |z_1| \, K(z)\, dz \;=\; \int_0^\infty r\, k(r)\, dr,
  \label{eq:interface-constant}
\end{equation}
where $k$ is the one-dimensional marginal of $K$. The second form is the flat
calculation: for an interface $\{x_1 = 0\}$, the pairs $(s, t)$ with
$s < 0 < t$ and $t - s = r$ have measure $r$, so the escaped heat per unit
length is $\int_0^\infty r\, k_\tau(r)\, dr = \sqrt\tau \int_0^\infty r\, k(r)\, dr$.
Summed over cells in \eqref{eq:heat-content}, every interface is counted from both
sides, so
\begin{equation}
  \Etau(\omega) \;\longrightarrow\; c_K \sum_k \operatorname{Per}(\Omega_k).
  \label{eq:gamma-limit}
\end{equation}
For the Gaussian, $k$ is $N(0, 2)$ and $c_G = \tfrac12 \sqrt{4/\pi} = 1/\sqrt\pi$.
This is the constant $c_0 = (|B^{d-1}|/|S^{d-1}|) \int |z|\, G(z)\, dz$ of
\cite[eq.~(A.2)]{esedoglu2015threshold}: in $d = 2$ it is $\pi^{-1}\sqrt\pi = 1/\sqrt\pi$,
and the same value follows from \eqref{eq:interface-constant} because
$\int |z_1| K = (|B^{d-1}|/|S^{d-1}|)\cdot 2 \int |z| K$ for a radial kernel. The
rigorous statement --- $\Gamma$-convergence of \eqref{eq:relaxed} in $L^1$ to
$c_K \sum_k \operatorname{Per}(\Omega_k)$, with compactness --- is
\cite[Proposition~A.1]{esedoglu2015threshold} for kernels that are smooth,
nonnegative, radial, radially non-increasing, of unit mass with a finite first
moment, and satisfy $|\nabla G(z)| \lesssim G(z/2)$; and
\cite[Theorem~5.1]{esedoglu2017kernels} for kernels that are merely
nonnegative, symmetric, in $L^1$ with $|z|K \in L^1$, and positive in a
neighbourhood of the origin. The kernel used here fails the first set
(it is singular at the origin, \S\ref{sec:kernel}) and satisfies the second.
Both theorems are proved on a periodic box in $\mathbb{R}^d$; neither covers a
curved surface or a finite-element discretisation. That transfer is assumed
here, not proved.

\section{The surface finite-element diffusion step}
\label{sec:diffusion}

\subsection{One backward-Euler step}

The semi-discrete P1 heat equation on the surface is $M \dot y + K y = 0$
\cite{dziuk2013surfacefem}. One backward-Euler step of length $\tau$ from
$y(0) = \chi_k$ is
\begin{equation}
  (M + \tau K)\, y_k \;=\; M \chi_k,
  \qquad\text{i.e.}\qquad
  y_k \;=\; \Atau \chi_k, \quad \Atau \;=\; (M + \tau K)^{-1} M .
  \label{eq:be-step}
\end{equation}
This is the only linear solve in the method. $M + \tau K$ is symmetric positive
definite, is factorised once per $(\text{level}, \tau)$ pass --- in the step
loop, again for the probe baseline and each probe factor, and for the final
full-budget assignment --- and the $N$ right-hand sides are back-substituted.

\begin{tcolorbox}[colback=blue!4!white, colframe=blue!40!black,
                  title={Code: the diffusion step}, fontupper=\small]
\codefile{src/partition/mbo\_auction.py}: \code{make\_diffusion\_solver(mesh, tau)}
returns \code{factorized((M + tau*K).tocsc())} (SuperLU);
\code{diffuse\_indicators(solve, mesh, labels, n\_cells)} forms
$M\chi$ sparsely and solves column by column, accumulating
$\sum_k (M\chi_k)^\top y_k$ for \S\ref{sec:instruments}.
\end{tcolorbox}

\subsection{The step is an exponential mixture of heat semigroups}
\label{sec:mixture}

Everything that is specific to this kernel follows from one identity. Let
$L = M^{-1} K$; it is diagonalisable with real eigenvalues
$0 = \lambda_1 \le \lambda_2 \le \cdots$ by the generalised eigenproblem
$K\phi = \lambda M\phi$, and $e^{-tL}$ is the semi-discrete heat semigroup.

\begin{prop}[Mixture representation]
\label{prop:mixture}
\begin{equation}
  \Atau \;=\; (I + \tau L)^{-1} \;=\; \int_0^\infty \frac{1}{\tau}\, e^{-t/\tau}\, e^{-tL}\, dt,
  \label{eq:mixture}
\end{equation}
i.e.\ one backward-Euler step is the heat semigroup averaged over an
exponentially distributed time $t \sim \mathrm{Exp}(\text{mean } \tau)$. The same
identity holds in the continuum, $(I - \tau\Delta)^{-1} = \int_0^\infty
\tau^{-1} e^{-t/\tau} e^{t\Delta}\, dt$.
\end{prop}
\begin{proof}
On each eigenvector, $\int_0^\infty \tau^{-1} e^{-t/\tau} e^{-t\lambda}\, dt
= \tau^{-1} (\tau^{-1} + \lambda)^{-1} = (1 + \tau\lambda)^{-1}$; the eigenvectors
span $\mathbb{R}^V$.
\end{proof}

\subsection{Structural properties}
\label{sec:structure}

\begin{prop}
\label{prop:structure}
\begin{enumerate}[label=(\roman*), itemsep=1pt]
  \item \textbf{Partition of unity.} $\sum_k y_k = \mathbf 1$ exactly, for every
    labelling.
  \item \textbf{Symmetry.} $M\Atau = M(M + \tau K)^{-1}M$ is symmetric positive
    definite, so $\chi \mapsto -\chi^\top M \Atau \chi$ is concave.
  \item \textbf{Positivity is not guaranteed on the mesh.} In the continuum
    $\Atau$ has a positive kernel (\S\ref{sec:kernel}). On a triangulation, $K$ has
    positive off-diagonal entries across obtuse angles and $M$ is not diagonal, so
    $e^{-tL}$ need not be entrywise nonnegative (on a $40 \times 36$ torus mesh its
    minimum entry is $-1.4 \times 10^{-2}$ at $t = 0.1 h^2$), and no argument
    gives nonnegativity of the mixture $\Atau$ either. In practice it holds: on
    every torus mesh tested at the production $\tau$ the minimum entry of $\Atau$
    is positive ($+2.7 \times 10^{-5}$ at $40 \times 36$, $N = 10$;
    $+8.7 \times 10^{-9}$ at $60 \times 52$, $N = 100$; $+3 \times 10^{-10}$ and
    $0$ to twelve digits at $100 \times 96$, $N = 100$ and $400$), so the caveat
    is inactive on the ladder. (\code{numerics.yaml}:
    \code{positivity\_and\_pairing.expm\_check}, \code{pairing\_cases[*].min\_entry\_A\_tau}.)
\end{enumerate}
\end{prop}
\begin{proof}
(i) $K\mathbf 1 = 0$ gives $(M + \tau K)\mathbf 1 = M\mathbf 1$, so
$\Atau \mathbf 1 = \mathbf 1$, and $\sum_k \chi_k = \mathbf 1$. (ii) $(M+\tau K)^{-1}$
is SPD and $M$ is symmetric, so $M(M+\tau K)^{-1}M$ is SPD. (iii) is the standard
failure of the discrete maximum principle for the cotangent Laplacian on
non-Delaunay meshes, and the numbers are direct computation. Nothing in the
method relies on positivity: the assignment is invariant under adding a common
constant to a vertex's scores, and the balanced threshold does not compare
against $\tfrac12$.
\end{proof}

\subsection{The continuum kernel and its interface constant}
\label{sec:kernel}

By \eqref{eq:mixture} the kernel of $(I - \tau\Delta)^{-1}$ on $\mathbb{R}^2$ is
\begin{equation}
  \Ktau(r) \;=\; \int_0^\infty \frac{1}{\tau} e^{-t/\tau}\, \frac{e^{-r^2/4t}}{4\pi t}\, dt
  \;=\; \frac{1}{2\pi\tau}\, K_0\!\Bigl(\frac{r}{\sqrt\tau}\Bigr),
  \label{eq:resolvent-kernel}
\end{equation}
with $K_0$ the modified Bessel function. From the mixture, without evaluating
\eqref{eq:resolvent-kernel}: $\Ktau$ is nonnegative, radial, of unit mass, has all
moments $\int |z|^p \Ktau = \mathbb{E}_t[\int |z|^p \Gt] \propto \mathbb{E}[t^{p/2}] < \infty$,
and is positive (indeed logarithmically singular) at the origin. It therefore
lies in the class of \cite[Theorem~5.1]{esedoglu2017kernels}; its Fourier symbol
$(1 + \tau|\xi|^2)^{-1}$ is positive, so it also satisfies the hypothesis of
Proposition~\ref{prop:dissipation}'s continuum version
\cite[\S5.2]{esedoglu2015threshold}. Its tails are exponential rather than
Gaussian: the one-dimensional marginal is
\begin{equation}
  k_\tau(s) \;=\; \frac{1}{2\sqrt\tau}\, e^{-|s|/\sqrt\tau},
  \label{eq:marginal}
\end{equation}
because restricting the symbol to the $\xi_1$-axis gives $(1 + \tau\xi_1^2)^{-1}$,
the transform of \eqref{eq:marginal}.

\begin{prop}[Interface constant of the backward-Euler kernel]
\label{prop:constant}
$c_{\Ktau} = \tfrac12$. Hence, for the kernel actually used,
\begin{equation}
  \Etau(\omega) \;\longrightarrow\; \frac12 \sum_k \operatorname{Per}(\Omega_k)
  \;=\; \text{total interface length, each interface counted once.}
  \label{eq:gamma-resolvent}
\end{equation}
\end{prop}
\begin{proof}
Either from \eqref{eq:marginal}:
$\int_0^\infty r\, k_1(r)\, dr = \tfrac12 \int_0^\infty r e^{-r}\, dr = \tfrac12$.
Or from the mixture: the Gaussian at time $t$ escapes $\sqrt{t/\pi}$ per unit
length, and $\mathbb{E}[\sqrt t\,] = \sqrt\tau\,\Gamma(\tfrac32) = \tfrac12\sqrt{\pi\tau}$,
so $\tau^{-1/2}\, \mathbb{E}[\sqrt{t/\pi}\,] = \tfrac12$.
\end{proof}

\begin{remark}[The factor of two in the perimeter convention]
Every perimeter this project reports --- \code{final\_perimeter},
\code{arm\_report.yaml}, the deliverable tables --- is $\sum_k \operatorname{Per}(\Omega_k)$,
which counts each interface twice. By \eqref{eq:gamma-resolvent}, $\Etau$
approximates \emph{half} of that quantity, i.e.\ the same thing
\code{label\_boundary\_length} measures (\S\ref{sec:instruments}). With the
Gaussian it would approximate $(2/\sqrt\pi) \approx 1.128$ times the interface
length instead.
\end{remark}

\subsection{The per-step displacement}
\label{sec:displacement}

\begin{prop}[Displacement per step]
\label{prop:displacement}
Let $\Sigma$ have a smooth boundary with (geodesic) curvature $\kappa$ at a point,
positive when $\Sigma$ is locally convex. One step \eqref{eq:mbo2} with the
Gaussian $G_\tau$ moves the boundary inward by $\tau\kappa$; with the
backward-Euler kernel $\Ktau$ it moves it by
\begin{equation}
  \delta \;=\; \frac{\tau\kappa}{2},
  \label{eq:displacement}
\end{equation}
both to leading order in $\tau$. One backward-Euler MBO step is therefore
mean-curvature flow over time $\tau/2$, not $\tau$.
\end{prop}
\begin{proof}
At signed distance $d$ outside $\Sigma$, the Taylor expansion of the diffused
indicator is \cite[\S5]{esedoglu2015threshold}
\begin{equation}
  (\Gt * \chi_\Sigma)(d) \;=\; \frac12 - \frac{d + t\kappa}{\sqrt{4\pi t}} + O(d^2, t),
  \label{eq:taylor}
\end{equation}
so the Gaussian threshold set is $d = -t\kappa$. Averaging \eqref{eq:taylor} over
$t \sim \mathrm{Exp}(\tau)$ and setting it to $\tfrac12$,
\[
  d\,\mathbb{E}[t^{-1/2}] + \kappa\, \mathbb{E}[t^{1/2}] = 0
  \;\Longrightarrow\;
  d = -\kappa\, \frac{\mathbb{E}[t^{1/2}]}{\mathbb{E}[t^{-1/2}]}
    = -\kappa\, \frac{\sqrt\tau\, \Gamma(3/2)}{\tau^{-1/2}\, \Gamma(1/2)}
    = -\frac{\kappa\tau}{2}. \qedhere
\]
\end{proof}

\begin{remark}[Second moments do not determine the speed]
Both kernels have second moment $2\tau$ per coordinate (expand
$(1 + \tau|\xi|^2)^{-1}$ and $e^{-\tau|\xi|^2}$ to first order), so one might
expect the same speed. The displacement is instead $\delta = \tfrac{\kappa}{2}\, \rho_K$ with
$\rho_K = \int y_2^2 K(0, y_2)\, dy_2 \big/ \int K(0, y_2)\, dy_2$ the ratio of the
transverse second moment to the marginal density \emph{on the interface}; $\rho_K$
is $2\tau$ for the Gaussian and $\tau$ for $\Ktau$ (a direct computation with
\eqref{eq:resolvent-kernel}, using $\int_0^\infty s^2 K_0(s)\, ds = \int_0^\infty K_0(s)\, ds = \pi/2$).
The mixture proof avoids the Bessel integrals.
Every identity in \S\ref{sec:mixture}--\S\ref{sec:displacement} --- the mixture
against the closed-form kernel, both interface constants, the displacement by
both routes, the Bessel integrals, the second moments, and the centre values of
\S\ref{sec:overmerge} --- is evaluated numerically in \code{numerics.yaml}:
\code{constants}.
\end{remark}

\subsection{Numerical check on the production mesh}
\label{sec:check}

Proposition~\ref{prop:displacement} is for the continuum kernel and to leading
order in $\tau$; the implementation uses \eqref{eq:be-step} on a triangulation.
For a disc the continuum problem is exactly solvable, which gives a reference
that is not leading-order. With $x = R/\sqrt\tau$ and $u = r/\sqrt\tau$, the
radial solution of $(1 - \tau\Delta)\, y = \chi_{B_R}$ is
\begin{equation}
  y(u) \;=\;
  \begin{cases}
    1 - x K_1(x)\, I_0(u), & u \le x,\\
    x I_1(x)\, K_0(u), & u \ge x,
  \end{cases}
  \label{eq:disc-exact}
\end{equation}
(continuous with continuous derivative at $u = x$ by the Wronskian
$I_0 K_1 + I_1 K_0 = 1/x$; the centre value is \eqref{eq:centre}), so the
$\tfrac12$-threshold radius $u^\ast$ and the exact area loss
$\pi\tau (x^2 - u^{\ast 2})$ follow by a one-dimensional root find. At the $x$ of
the table below it is $1.09$--$1.16 \times \pi\tau$: the leading-order $\pi\tau$
is low by $9$--$16\%$ at these $x$, because the resolvent's exponential tails
make its $O(1/x)$ correction about four times the Gaussian's (whose exact loss
is $1.02$--$1.04 \times 2\pi\tau$ at the same $x$).

On the $348 \times 328$ torus mesh of the deliverables ($V = 114{,}144$,
$\hmean = 0.0171$, $\hmax = 0.0311$), take the set of vertices within
\emph{ambient} distance $R_d$ of a point on the outer equator --- not a geodesic
disc; its lumped area differs from $\pi R_d^2$ by $-1.4$ to $+0.8\%$ --- apply one
step of \eqref{eq:be-step} and a $\tfrac12$-threshold, and measure the lumped
area change. The result depends on where the centre sits relative to the mesh:
the boundary ring of vertices flips in groups, so a single placement carries
granularity noise. The table gives the vertex-centred placement $(1.6, 0, 0)$
and the mean over 24 centres drawn uniformly within one grid cell around it, as
ratios to the exact continuum loss \eqref{eq:disc-exact} (all values:
\code{numerics.yaml}: \code{disc\_check}; seed and draw order in the provenance
block):

\begin{center}\footnotesize
\begin{tabular}{@{}rrrrrrrr@{}}
\toprule
$R_d$ & $\tau$ & $\sqrt\tau/\hmean$ & $x$ & exact loss & vertex-centred & jittered mean $\pm$ std & vs.\ Gaussian exact \\
\midrule
0.20 & 0.004 & 3.69 & 3.16 & $-0.01454$ & 0.90 & $0.97 \pm 0.06$ & 0.54 \\
0.25 & 0.004 & 3.69 & 3.95 & $-0.01392$ & 0.89 & $0.96 \pm 0.08$ & 0.52 \\
0.30 & 0.006 & 4.52 & 3.87 & $-0.02095$ & 0.84 & $0.93 \pm 0.04$ & 0.50 \\
0.20 & 0.002 & 2.61 & 4.47 & $-0.00684$ & 0.77 & $0.91 \pm 0.09$ & 0.48 \\
\bottomrule
\end{tabular}
\end{center}

Averaged over placement the mesh reaches $0.91$--$0.97$ of the exact continuum
loss and $0.48$--$0.54$ of the exact Gaussian value; the vertex-centred
placement is the low tail of that distribution (an earlier version of this
section reported it alone and read a $10$--$23\%$ deficit into it). Of the
remaining $3$--$9\%$ shortfall, surface curvature accounts for $1.4$--$3.1\%$
(the Gaussian curvature at the outer equator is $1/(r(R + r)) = 1.042$, which
lowers the boundary's geodesic curvature below $1/R_d$ by $K R_d^2/3$, i.e.\
$3.1\%$ at $R_d = 0.3$); the rest is finite-element and threshold-granularity
error at a per-step displacement of only $0.25$--$0.63\, \hmean$ (implied by the
measured vertex-centred loss, \code{displacement\_over\_h\_mean}). What the check establishes is that
the FEM step moves an interface at \emph{half} the Gaussian rate to within the
granularity of the test, as Proposition~\ref{prop:displacement} says. This is a
one-off check of the derivation, not a measured study.

\section{The balanced threshold}
\label{sec:assignment}

\subsection{The constrained linearisation is a transportation problem}

Minimising the linearisation \eqref{eq:linearisation} under the equal-area
constraint, discretised with a weight matrix $W$ for the inner product, is
\begin{equation}
  \max_{\chi'} \;\sum_k \ip{\chi'_k}{W y_k}
  \quad\text{s.t.}\quad
  \sum_k \chi'_{ik} = 1 \;\;\forall i, \qquad
  \sum_i v_i\, \chi'_{ik} = \Atgt \;\;\forall k, \qquad
  \chi'_{ik} \in \{0, 1\}.
  \label{eq:assignment}
\end{equation}
The implementation takes $W = \Dlump$, so the objective is
$\sum_{i} v_i\, y_{i,\omega'(i)}$ and \eqref{eq:assignment} is the transportation
problem (LP) of \codefile{docs/math/09-balanced-readout/} \S6.2 with score matrix
$s_{ik} = y_{ik}$: supplies $v_i$ at the vertices, demands $\Atgt$ at the cells.
With unit vertex masses it would be the assignment problem of
\cite{jacobs2018auction}; with the real lumped masses it is a transportation
problem \cite{ahuja1993network}, which is why exact balance is generically
unattainable (doc~09, \S6.5) and why a granularity bar exists at all.

\subsection{Dual, shifted argmax, and the two-phase antecedent}
\label{sec:dual}

Doc~09 derives the dual of the LP relaxation and shows (its Proposition~6.1) that
complementary slackness forces the assignment onto
\begin{equation}
  \omega'(i) \;=\; \argmax_k \bigl( y_{ik} + \psi_k \bigr),
  \label{eq:shifted-argmax}
\end{equation}
in the code's sign convention, where raising $\psi_k$ grows cell $k$, and that the
dual subgradient is $T_k - \Atgt$: grow the cells that are too small. The
$\psi_k$ are the Lagrange multipliers $\lambda^*$ of \cite[eq.~(11)]{jacobs2018auction}
and the \emph{prices} of the auction \cite{bertsekas1988auction}. For $N = 2$,
using $y_{i1} + y_{i2} = 1$,
\begin{equation}
  y_{i1} + \psi_1 \ge y_{i2} + \psi_2
  \;\Longleftrightarrow\;
  y_{i1} \;\ge\; \frac12 + \frac{\psi_2 - \psi_1}{2},
  \label{eq:ruuth-wetton}
\end{equation}
which is the threshold level shifted away from $\tfrac12$ until the volume is met
--- the scheme of \cite{ruuth2003volume}. \eqref{eq:shifted-argmax} is its
$N$-phase form.

\subsection{The inexact solver as implemented}
\label{sec:inexact}

\code{solve\_dual\_offsets} is a subgradient method on the piecewise-linear dual
with a decaying step --- descent on $D(\psi)$ in doc~09's convention, ascent on
the concave form in the code's --- returning the \emph{best} iterate visited
(doc~09, \S6.2--6.3). Three things are specific to its use here:
\begin{enumerate}[itemsep=1pt]
  \item \textbf{Scale.} The step schedule was calibrated on $\log$-density
    scores. Diffused indicators live on $[0, 1]$, so with
    \code{normalize\_scores=True} the scores are divided by the median gap
    between each vertex's top two scores (the margin an offset must cover to
    flip a vertex), the schedule \code{normalized\_eta0 = 2.0},
    \code{dual\_decay = 0.02}, \code{normalized\_dual\_iters = 2000} is run on the
    normalised scores, and $\psi$ is rescaled to the caller's units. Correctness
    of \eqref{eq:shifted-argmax} is scale-free; only convergence is not.
  \item \textbf{Quality bar.} A solve is accepted at
    $\max_k |T_k - \Atgt|/\Atgt \le \text{bar}$ with
    $\text{bar} = \max(2\,\gamma, \text{A-parity})$, where
    $\gamma = \max_i v_i / \Atgt$ is the one-vertex granularity and A-parity is
    approach~A's own stall level at equal budget (report~07). A-parity is a
    two-entry lookup, for $(V, N) = (47{,}488, 300)$ and $(114{,}144, 300)$ only;
    for every other configuration --- all $N = 100, 400$--$1000$ runs and every
    non-torus surface --- the bar is exactly $2\gamma$. It is a solver-failure
    detector, not a validity test: equal area is what the step optimises, so
    hitting it says the dual converged.
  \item \textbf{Early stopping is production-only.} In-loop solves stop at the
    bar; every level ends with a full-budget re-assignment, and in-loop area
    figures are tagged \code{early\_stop\_capped} so they are never quoted as
    assignment quality (report~07, artefact~5).
\end{enumerate}

\begin{tcolorbox}[colback=blue!4!white, colframe=blue!40!black,
                  title={Code: the balanced threshold}, fontupper=\small]
\codefile{src/partition/mbo\_auction.py}: \code{balanced\_assign(scores, lumped\_mass, target, bar, early\_stop)}
$\to$ \codefile{src/partition/balanced\_readout.py}: \code{solve\_dual\_offsets}
with \code{BalancedReadoutConfig(normalize\_scores=True, dual\_early\_stop\_rel=bar or None)};
\code{assignment\_margin\_scale} is the median top-two gap;
\codefile{src/partition/arm\_harness.py}: \code{assignment\_quality\_bar},
\code{vertex\_granularity}.
\end{tcolorbox}

\begin{remark}[Why an inexact standard solver, and what may be claimed for it]
\label{rem:solver-choice}
Nothing in \S\ref{sec:assignment} is a new method for the transportation dual.
Subgradient ascent on that dual --- grow the cells that are too small --- is
the standard solver for this class of problem
\cite{aurenhammer1998minkowski, merigot2011multiscale, peyre2019computational},
and doc~09 (its Proposition~6.1) proves that the implemented update \emph{is}
that subgradient. What is specific to this work is the \emph{decision} to use it
inside the threshold step in place of the exact auction of
\cite{jacobs2018auction}, and the reasons are these, each on record:
\begin{enumerate}[itemsep=1pt]
  \item the adaptation was forced by the mesh regardless of solver --- the
    auction of \cite{jacobs2018auction} is stated for unit-volume points, and
    with real lumped masses the problem is a transportation problem whose
    optimum is generically fractional (\S\ref{sec:assignment}, doc~09 \S6.5), so
    a granularity bar is the correct stopping rule for \emph{any} solver;
  \item a validated primitive already existed: the routine was built for the
    balanced readout and qualified through the Phase~0 harness
    (\codefile{docs/experiments/07-phase0-shared-harness/}), which reproduces the
    reference Phase~2 perimeter bitwise;
  \item inexactness is safe \emph{because the scheme is iterated}: report~07
    records that the solver looked inadequate only on a cold-start fixture that
    no iterated method faces, and that from the first step every assignment
    starts from an already-balanced partition and converges under the bar ---
    a conditional judgment, checked rather than assumed;
  \item the cost is bounded, not waived: the dissipation theorem is lost
    (\S\ref{sec:outside}), the inequality that replaces it holds for any priced
    labelling (Proposition~\ref{prop:survives}) and its slack is measured on
    every run (\S\ref{sec:survives}).
\end{enumerate}
The limit of the claim: the choice \emph{sufficed}, and its price is known. It
is \textbf{not} established that it is better, faster or more accurate than the
auction, because no auction was implemented or measured on this problem; the
one alternative that was measured, entropic (Sinkhorn) regularisation
\cite{cuturi2013sinkhorn}, lost to the incumbent on five of six score matrices
(report~07). This remark exists so that the mathematical record carries the
reasoning behind the choice and the disclaimer with it, not only the
consequences.
\end{remark}

\section{Energy dissipation: what transfers and what does not}
\label{sec:monotonicity}

\subsection{The theorem, for the consistent pairing with an exact assignment}
\label{sec:theorem}

Discretise \eqref{eq:relaxed} with the consistent mass matrix:
\begin{equation}
  E^{M}_\tau(\chi) \;=\; \frac{1}{\sqrt\tau}\Bigl[\, |S| - \sum_k \chi_k^\top M \Atau \chi_k \Bigr]
  \;=\; \frac{1}{\sqrt\tau}\Bigl[\, |S| - \sum_k \ip{\chi_k}{M y_k} \Bigr].
  \label{eq:E-consistent}
\end{equation}
By Proposition~\ref{prop:structure}(ii) this is concave on the convex hull of
labellings, and its linearisation at $\chi$ is
$-\tfrac{2}{\sqrt\tau} \sum_k \ip{\chi'_k}{M y_k}$. So if the assignment step
\emph{exactly} maximised $\sum_k \ip{\chi'_k}{M y_k}$ over labellings with equal
discrete area, Proposition~\ref{prop:dissipation} would give
$E^M_\tau(\chi') \le E^M_\tau(\chi)$ at every step, constraint included --- the
discrete form of the auction-dynamics theorem of \cite{jacobs2018auction}.

\subsection{Why the implementation is outside it}
\label{sec:outside}

Two hypotheses fail, independently.
\begin{enumerate}[itemsep=1pt]
  \item \textbf{The pairing.} The assignment maximises
    $\sum_k \ip{\chi'_k}{\Dlump y_k} = \sum_k \chi_k'^\top \Dlump\Atau \chi_k$ with the
    \emph{lumped} $\Dlump$, while $y_k$ was produced with the \emph{consistent}
    $M$, and $\Dlump\Atau = \Dlump (M + \tau K)^{-1} M$ is not symmetric. Write
    $\Dlump\Atau = S + Z$ with $S$ symmetric and $Z$ skew;
    $\|Z\|_F / \|\Dlump\Atau\|_F$ is $2 \times 10^{-3}$ on the $40 \times 36$ and
    $60 \times 52$ torus meshes and $0.9$--$1.2 \times 10^{-3}$ on the
    $100 \times 96$ production base. The lumped energy
    $E^{\Dlump}_\tau = \tau^{-1/2}[\,|S| - \sum_k \chi_k^\top S \chi_k]$ (the skew part
    drops out of a quadratic form) is concave iff $S \succeq 0$, which is not
    guaranteed but was true on every mesh tested at the production $\tau$:
    smallest eigenvalue $+4 \times 10^{-6}$ ($40 \times 36$, $N = 10$;
    $60 \times 52$, $N = 100$) and $+6 \times 10^{-7}$, $+2 \times 10^{-6}$
    ($100 \times 96$ at $N = 100$, $400$), against largest eigenvalues of
    $3 \times 10^{-3}$ to $2 \times 10^{-2}$. The operative failure is
    elsewhere: the linearisation of $-\chi^\top S \chi$ at $\chi$ is $-2 S\chi$,
    whereas the assignment maximises $\ip{\chi'}{(S + Z)\chi}$. The step is
    therefore the exact minimise-the-linearisation step of
    Proposition~\ref{prop:dissipation} for $E^{\Dlump}_\tau$ \emph{perturbed by
    the skew term} $\ip{\chi'}{Z\chi}$. The matrix-norm figure above is an upper
    bound on that perturbation, not its size: at an actual pair of labellings
    (the geodesic balanced init at the project seed, then one full-budget
    balanced step) $\ip{\chi'}{Z\chi}/\ip{\chi'}{S\chi}$ is $1.1 \times 10^{-5}$
    on $40 \times 36$ and $2.0 \times 10^{-5}$ on $60 \times 52$, and the skew
    term changes the unconstrained argmax at one vertex of $1{,}440$ and one of
    $3{,}120$ respectively (\code{numerics.yaml}:
    \code{pairing\_cases[*].labelling\_pair}). All eigenvalue and asymmetry
    figures in this item: \code{pairing\_cases[*]}. Monotonicity of
    $E^{\Dlump}_\tau$ can fail only through that term and through inexactness
    (next item); tracking $E^M_\tau$ instead is simply the wrong pairing, since
    its gradient $-2 M\Atau\chi$ is not what the assignment sees at all.
  \item \textbf{Exactness.} \cite{jacobs2018auction} obtain a theorem because
    the auction solves the assignment exactly (to $\varepsilon$-complementary
    slackness). \code{solve\_dual\_offsets} returns the best iterate of a
    subgradient method with a residual area error, and the early-stopped in-loop
    solves are further from optimal by design.
\end{enumerate}
Report~08 records the consequence: $E^{\Dlump}_\tau$ \emph{increases} on 4 of
96 active steps at $N = 300$, consistent with two small perturbations of an
otherwise monotone scheme. Neither energy is therefore a gate. The discrete
forms of the dissipation statement exist in the literature and are the right
reference here: \cite[\S4.4]{esedoglu2017kernels} state that two-phase graph MBO
with a symmetric weight matrix $W$ decreases its energy at every step when $W$ is
positive semi-definite, and exhibit a $W = A + \lambda I$ ($\lambda < 2$) for which
it cycles between two configurations (their Example~3, Fig.~3); their numbered
Propositions~4.7--4.9 concern modified algorithms with split weights and are
not the statement used here. For the constrained multiclass scheme
\cite[\S5.2]{jacobs2018auction} state the same condition (``as long as $W$ is
positive semi-definite, the graph heat content is concave''), which is
Proposition~\ref{prop:dissipation} applied to a graph energy.

\subsection{The inequality that survives}
\label{sec:survives}

\begin{prop}[Descent inequality]
\label{prop:survives}
Let $\omega'$ be \emph{any} labelling of the form \eqref{eq:shifted-argmax} for
some $\psi \in \mathbb{R}^N$ --- exact or not --- and $\omega$ the previous
labelling. Then
\begin{equation}
  \sum_k \ip{\chi'_k}{\Dlump y_k} \;-\; \sum_k \ip{\chi_k}{\Dlump y_k}
  \;\ge\; \sum_k \psi_k \bigl( T_k - T'_k \bigr),
  \label{eq:survives}
\end{equation}
where $T_k, T'_k$ are the old and new discrete areas.
\end{prop}
\begin{proof}
For every vertex $i$, the argmax gives
$y_{i\omega'(i)} + \psi_{\omega'(i)} \ge y_{i\omega(i)} + \psi_{\omega(i)}$.
Multiply by $v_i > 0$ and sum: the $y$ terms give the left side of
\eqref{eq:survives}, and $\sum_i v_i (\psi_{\omega(i)} - \psi_{\omega'(i)})
= \sum_k \psi_k (T_k - T'_k)$.
\end{proof}

The right-hand side is exactly the price of inexactness: at perfect balance
$T_k = T'_k = \Atgt$ it vanishes and \eqref{eq:survives} becomes the textbook
monotone step for the lumped linearisation; otherwise it scales like
$\|\psi\| \times$ (area error): the per-level maxima recorded in report~08's
\code{data.yaml} run from $1.3 \times 10^{-4}$ ($N = 100$) to $2.8 \times 10^{-2}$
($N = 1000$ on $V = 114{,}144$, level~1), relative to $|v^\top y_{\omega}|$. Because
\eqref{eq:survives} needs neither concavity nor exactness, a violation beyond
floating point can only mean the returned labels are not the argmax at the
returned $\psi$, or a unit mismatch between them; that makes it worth gating on,
and the measured violation is $\le 6.1 \times 10^{-8}$.

\begin{tcolorbox}[colback=blue!4!white, colframe=blue!40!black,
                  title={Code: energies and the gate}, fontupper=\small]
\codefile{src/partition/mbo\_auction.py}: \code{lyapunov\_energy(y, labels, v, total\_area, tau, inner\_consistent)}
returns \code{E\_lumped} $= E^{\Dlump}_\tau$ and \code{E\_consistent} $= E^M_\tau$;
\code{descent\_slack(y, v, psi, labels\_old, labels\_new, n\_cells)} returns
\code{lhs}, \code{rhs} of \eqref{eq:survives} and \code{violation\_rel} $=$ (rhs$-$lhs)/scale.
\codefile{testing/test\_mbo\_auction.py} gate G4 asserts \eqref{eq:survives}.
\end{tcolorbox}

\subsection{A pairing that would restore the theorem (not implemented)}
\label{sec:lumped-pairing}

\begin{remark}[Lumped--lumped pairing]
Replace \eqref{eq:be-step} by the mass-lumped step
$(\Dlump + \tau K)\, y_k = \Dlump \chi_k$, i.e.\ $\Atau^{\Dlump} = (\Dlump + \tau K)^{-1}\Dlump$,
and keep the assignment as it is. Then $\Dlump\Atau^{\Dlump} = \Dlump(\Dlump + \tau K)^{-1}\Dlump$
is symmetric positive definite, the lumped energy $E^{\Dlump}_\tau$ is concave,
its linearisation is the objective the assignment already maximises, and with an
exact assignment Proposition~\ref{prop:dissipation} applies verbatim. Mass
lumping is a standard, first-order-consistent variant of the P1 heat step
\cite{dziuk2013surfacefem}; Propositions~\ref{prop:mixture}--\ref{prop:displacement}
hold unchanged with $L = \Dlump^{-1}K$. This changes the diffusion operator and
therefore every measured number, so it is a separate experiment with its own
falsifiers; it is recorded here because it is the one-line change that removes the skew
term of \S\ref{sec:outside}(1) exactly --- with $W = \Dlump(\Dlump + \tau K)^{-1}\Dlump$,
which is symmetric positive definite, the lumped energy is the graph heat
content of \cite[\S5.2]{jacobs2018auction} and the assignment objective is its
linearisation, so Proposition~\ref{prop:dissipation} applies --- leaving only
the inexactness of \S\ref{sec:outside}(2). Since the skew term is
$\lesssim 2 \times 10^{-3}$ in matrix norm and $\sim 10^{-5}$ at an actual
labelling pair, the experiment would be isolating a small effect.
\end{remark}

\section{The time step: a two-sided window}
\label{sec:tau}

\subsection{The rule}

Per mesh level,
\begin{equation}
  \tau \;=\; \min\bigl( (c\, \hmean)^2,\; (\rho\, \Rcell)^2 \bigr),
  \qquad c = 4, \quad \rho = 1,
  \label{eq:tau-rule}
\end{equation}
so $\sqrt\tau$ is $c$ mean edges unless that would exceed the cell radius. The
two branches are the two sides of the window; the constants are pre-registered
and the measured band that passes every gate is in report~08 and
\codefile{CLAUDE.md}.

\subsection{Freeze}
\label{sec:freeze}

A vertex's label can change only if the diffused field crosses its neighbour's
at that vertex, which requires the continuum interface to move at least a mesh
spacing. With Proposition~\ref{prop:displacement} and $\kappa \sim 1/\Rcell$ for a
cell-scale feature,
\begin{equation}
  \delta = \frac{\tau}{2\Rcell} \;\gtrsim\; h
  \;\Longleftrightarrow\;
  \tau \gtrsim 2\, h\, \Rcell
  \;\Longleftrightarrow\;
  c \gtrsim \sqrt{2\Rcell/h} \quad\text{on the uncapped branch.}
  \label{eq:freeze}
\end{equation}
The implementation reports the same criterion without the factor $2$ (it was
written with the Gaussian displacement $\tau\kappa$):
\code{c\_lo\_hmean} $= \sqrt{\Rcell/\hmean}$ and
\code{c\_lo\_hmax\_local} $= \sqrt{\Rcell/\hmax}$. The factor shifts the threshold
by $\sqrt2$, inside the margin the constants were chosen with (the frozen
$c = 2$ / converged $c = 4$ pair is the NC3-CAL control of
\codefile{testing/test\_mbo\_auction.py}; report~08 \S3 cites
\cite{vangennip2014mean} for the freeze regime); it is noted, not corrected, because the ratio is a
diagnostic and every measured run quotes the reported value. Both are
order-of-magnitude statements: the expansion behind
Proposition~\ref{prop:displacement} assumes $h \ll \sqrt\tau \ll \Rcell$, and at
the coarsest levels neither inequality is strong. The rigorous statement in this
direction is for graphs: \cite[Theorem~4.2]{vangennip2014mean} proves the
two-phase graph scheme is stationary for every initial set when
$\tau < \log(3/2)/\rho(\Delta) \approx 0.4/\rho(\Delta)$, $\rho$ the spectral
radius. For a mesh Laplacian $\rho(L) \sim C/h^2$, so that bound is
$\sqrt\tau \lesssim 0.6\,h/\sqrt C$ --- a much weaker (necessary-for-motion)
condition than \eqref{eq:freeze}, which involves $\Rcell$.

\subsection{Over-merge}
\label{sec:overmerge}

On the other side, when $\sqrt\tau$ approaches $\Rcell$ a cell's own diffused
indicator no longer dominates at its own centre. From the mixture, for a disc of
radius $R$ the diffused indicator at the centre is
\begin{equation}
  y(0) \;=\; \int_0^\infty \frac{e^{-t/\tau}}{\tau}\bigl(1 - e^{-R^2/4t}\bigr)\, dt
  \;=\; 1 - x\, K_1(x), \qquad x = \frac{R}{\sqrt\tau},
  \label{eq:centre}
\end{equation}
($K_1$ the modified Bessel function; the Gaussian value is $1 - e^{-x^2/4}$).
At the cap $\rho = 1$, $x = 1$ and $y(0) = 1 - K_1(1) = 0.398$: the cell's own
signal at its core is below $\tfrac12$, i.e.\ the \emph{sum} of its neighbours'
diffused indicators already exceeds it there. The argmax compares against the
\emph{largest} neighbour, so this does not by itself imply core loss; it marks
the scale at which core loss becomes possible. At $x = 2$ and $4$, $y(0) = 0.720$ and $0.950$.
This is why $\rho = 1$ is called the \emph{loosest defensible} cap: it is the
point at which the diffusion has spread over the whole cell, and anything
tighter is a preference. Assignment quality cannot see over-merging --- the
balanced threshold hits its area target by construction whatever the geometry ---
so the over-merge instruments are geometric (\S\ref{sec:instruments}).

\subsection{The reported ratios, and anisotropy}
\label{sec:ratios}

\code{tau\_diagnostics} reports $\sqrt\tau/\hmean$, $\sqrt\tau/\hmax$,
$\sqrt\tau/\Rcell$, the two freeze constants, \code{cap\_active}, and the
\emph{freeze ratio}
\begin{equation}
  \frac{\sqrt\tau/\hmax}{\sqrt{\Rcell/\hmax}} \;=\; \sqrt{\frac{\tau}{\Rcell\, \hmax}}
  \;=\;
  \begin{cases}
    c\,\hmean/\sqrt{\Rcell \hmax} & \text{uncapped},\\[2pt]
    \sqrt{\Rcell/\hmax} & \text{capped},
  \end{cases}
  \label{eq:freeze-ratio}
\end{equation}
which is the quantity \codefile{CLAUDE.md} calls the $\tau$ freeze ratio. Two
consequences follow from \eqref{eq:freeze-ratio} and are borne out by the runs.
On the capped branch the ratio \emph{falls} as cells shrink relative to the mesh
($\Rcell \propto N^{-1/2}$), which is why the standard $100 \times 96$ base level
drops below $1$ at $N = 750$ and $1000$; on the uncapped branch it falls as cells
\emph{grow} relative to the mesh, which is why a well-resolved low-$N$
configuration runs at $0.5$--$0.6$ and passes every gate --- the ratio is not a
defect indicator on its own, and the measured band is the reference. The mesh is
anisotropic ($\hmax/\hmean = 1.81$ on the torus ladder), so a mean-based figure
flatters the coarse band: the rule is on $\hmean$ and the freeze ratio is on
$\hmax$, deliberately.

\subsection{The continuation probe}
\label{sec:probe}

$\Etau$ depends on the labels alone once $\tau$ is fixed, so a level that froze
and a level that converged both show a flat energy tail. The probe
\code{tau\_continuation\_probe} re-runs one balanced threshold with
$\tau' = f\tau$, $f \in \{2, 4\}$, and evaluates the resulting labels at the
\emph{level's} $\tau$: if a wider blur reaches a labelling with materially lower
$\Etau$ and moves a material fraction of vertices, the level was pinned. Because
a wider blur always re-tiebreaks some boundary vertices and shaves some energy,
the verdict needs calibrated thresholds (\code{nc3\_K}, \code{nc3\_f\_min}), and
until they are pinned by the NC3-CAL control the probe records and reports
\code{pinned = None} rather than a verdict. A pinned level would be annealed
($\tau \leftarrow 1.5\tau$, at most three times) and re-run. With the shipped
defaults \code{nc3\_K = nc3\_f\_min = None} the verdict is always \code{None}, so
the anneal branch is unreachable and every measured run reports
\code{n\_anneals = 0} vacuously --- the probe records, it does not yet decide.

\begin{tcolorbox}[colback=blue!4!white, colframe=blue!40!black,
                  title={Code: the window}, fontupper=\small]
\codefile{src/partition/mbo\_auction.py}: \code{tau\_diagnostics(mesh, n\_cells, cfg, tau\_override)}
implements \eqref{eq:tau-rule} and reports every ratio of \S\ref{sec:ratios};
\code{tau\_continuation\_probe} is \S\ref{sec:probe}; \code{MBOConfig.tau\_c},
\code{rho}, \code{anneal\_factor}, \code{max\_anneals}, \code{probe\_factors},
\code{nc3\_K}, \code{nc3\_f\_min}. Only \code{tau\_c}, \code{rho} and
\code{max\_iters} are config-reachable (\code{mbo:} block); the rest define the
measurement protocol.
\end{tcolorbox}

\section{Initialisation}
\label{sec:init}

The first labelling is one balanced assignment of geodesic scores:
\begin{enumerate}[itemsep=1pt]
  \item $N$ sites by greedy farthest-point sampling in the \emph{ambient}
    Euclidean metric of $\mathbb{R}^3$ (first site from the seeded RNG, each
    next site the vertex farthest from the current set) ---
    \code{farthest\_point\_sampling}, shared with the PGD seeded init so both
    start from the same site layout;
  \item $d_{ik}$ = shortest-path distance from site $k$ to vertex $i$ on the
    edge graph of the mesh, edge weights the Euclidean edge lengths, by $N$
    single-source Dijkstra searches returning the full $V \times N$ matrix
    (\code{dijkstra} with \code{indices=sites}) --- a graph metric that \emph{over}-estimates the
    geodesic distance (paths are confined to edges) but respects the surface
    topology, since it is built from \code{faces}; the heat method of
    \cite{crane2013geodesics} is not what is implemented;
  \item scores $s_{ik} = -d_{ik}^2$ and one full-budget balanced assignment.
\end{enumerate}
With $-d^2$ scores, \eqref{eq:shifted-argmax} is a \emph{power diagram} of the
graph metric with weights $\psi$ (doc~09, \S8; \cite{aurenhammer1998minkowski}):
the capacity-constrained Voronoi diagram that \cite{bogosel2021longest} also use
as an initialisation, computed here by the transportation dual rather than by
their gradient method, and on a surface. It is also, by construction, step~0 of
approach~C (capacity-constrained geodesic Lloyd), which is what makes the
attribution control of report~08 possible: running the init alone on the finest
mesh isolates what the dynamics add.

The balanced start is a \textbf{hard requirement}, measured rather than argued:
from a cold unbalanced Voronoi labelling the first balanced threshold leaves a
worst cell at $68.16\%$ area deviation, from the balanced one at $1.97\%$
(\codefile{docs/plans/PHASE1\_BC\_REPLACEMENT\_PLAN.md}, Phase~A, and the module
docstring; report~08 does not carry these two numbers).

\begin{tcolorbox}[colback=blue!4!white, colframe=blue!40!black,
                  title={Code: initialisation}, fontupper=\small]
\codefile{src/partition/mbo\_auction.py}: \code{geodesic\_balanced\_init(mesh, n\_cells, seed)};
\code{edge\_graph(mesh)}; \codefile{src/optimization/initialization.py}:
\code{farthest\_point\_sampling(vertices, n\_seeds, seed)};
\code{scipy.sparse.csgraph.dijkstra}.
\end{tcolorbox}

\section{The mesh ladder}
\label{sec:ladder}

Levels are the same mesh ladder Phase~1 uses. Within a level, steps repeat until
the fraction of vertices that change label is below \code{churn\_tol} $= 10^{-4}$
for \code{patience} $= 3$ consecutive steps, or \code{max\_iters} $= 200$ is hit
(never, in the measured runs: the most any level took is 66). Between levels the
labelling is carried by \emph{nearest embedded vertex}: each fine vertex takes
the label of the nearest coarse vertex in $\mathbb{R}^3$ (\code{cKDTree}).
Labels are interpolated rather than densities, because a hard partition has no
density to interpolate and the nearest-vertex map of a labelling is itself a
labelling. What the transfer preserves: each cell's territory up to a band of
width $O(h_{\text{coarse}})$ along its boundary, hence its area up to
$O(h_{\text{coarse}} \cdot \operatorname{Per}(\Omega_k))$ --- the first balanced
threshold on the fine level removes that error. What it does not preserve:
connectivity is not guaranteed across the transfer (a thin neck can be cut), and
$\tau$ is recomputed from \eqref{eq:tau-rule} on the new mesh, so the effective
diffusion length shrinks with $h$ on the uncapped branch. Report~08 records
fragmented cells born at one level and healed at the next; the ladder is what
heals them.

\begin{tcolorbox}[colback=blue!4!white, colframe=blue!40!black,
                  title={Code: ladder}, fontupper=\small]
\codefile{src/partition/mbo\_auction.py}: \code{run\_mbo\_ladder(meshes, n\_cells, cfg, init\_only)},
\code{mbo\_level(mesh, labels, n\_cells, cfg, level, run\_probe)},
\code{\_run\_steps}, \code{interpolate\_labels(old\_vertices, new\_vertices, labels)}.
\end{tcolorbox}

\section{Instruments}
\label{sec:instruments}

None of these is a gate on the dynamics; they are what the reports quote.
\begin{itemize}[itemsep=1pt]
  \item \textbf{Energies.} $E^{\Dlump}_\tau$ (\code{E\_lumped}, primary because
    its linearisation differs from the assignment objective only by the skew
    term of \S\ref{sec:outside}) and $E^M_\tau$ (\code{E\_consistent}),
    both by \eqref{eq:E-consistent} with the respective weight, and both
    $\to \tfrac12 \sum_k \operatorname{Per}(\Omega_k)$ by
    Proposition~\ref{prop:constant}. Descriptive trends only (\S\ref{sec:outside}).
  \item \textbf{Label boundary length.} The midpoint-cut construction: a
    triangle with two labels contributes the segment joining the midpoints of
    its two bi-labelled edges, one with three labels contributes three segments
    from those midpoints to the centroid (\code{label\_boundary\_length},
    \codefile{src/partition/balanced\_readout.py}). It counts each interface
    once; the Phase~2 perimeter counts each twice, and the measured ratio between
    them on an unoptimised labelling is $1.999$ (\codefile{CLAUDE.md}).
  \item \textbf{Compactness.} $Q_k = P_k^2 / (4\pi A_k)$ with $P_k$ the
    per-cell boundary length (same construction, each segment credited to both
    cells) and $A_k = T_k$; $Q = 1$ is a disc. Over-merging raises $Q$.
  \item \textbf{Core loss.} The number of cells whose own diffused-score peak
    vertex, $\argmax_i y_{ik}$, is not in their territory: the most direct
    over-merge symptom, since it says the cell's signal was swamped before the
    threshold was taken.
  \item \textbf{The three Phase~1 gates}, run on every level's labels through
    the same harness path as the final labels (\code{run\_gates}). For one-hot
    labels the dormant gate is vacuous and the area gate is the solver check of
    \S\ref{sec:inexact}; only connectivity has content, because nothing in
    \S\ref{sec:assignment} enforces it.
\end{itemize}

\begin{tcolorbox}[colback=blue!4!white, colframe=blue!40!black,
                  title={Code: instruments}, fontupper=\small]
\codefile{src/partition/mbo\_auction.py}: \code{lyapunov\_energy},
\code{per\_cell\_boundary\_lengths}, \code{compactness}, \code{core\_loss};
\codefile{src/partition/balanced\_readout.py}: \code{label\_boundary\_length};
\codefile{src/partition/arm\_harness.py}: \code{run\_gates}, \code{labels\_to\_one\_hot}.
\end{tcolorbox}

\section*{Summary table}

\begin{center}\small
\begin{tabular}{@{}p{5.0cm}p{4.3cm}p{5.8cm}@{}}
\toprule
Quantity & Status & Implementation \\
\midrule
Diffusion step $(M + \tau K) y_k = M\chi_k$ & derived (backward Euler of the P1 heat equation) & \code{make\_diffusion\_solver}, \code{diffuse\_indicators} \\
Mixture representation \eqref{eq:mixture} & proved & --- (analysis) \\
$\sum_k y_k = \mathbf 1$; $M\Atau$ SPD; no discrete positivity & proved & --- \\
Continuum kernel $\tfrac{1}{2\pi\tau}K_0(r/\sqrt\tau)$, marginal \eqref{eq:marginal} & derived & --- \\
Interface constant $c_{\Ktau} = \tfrac12$ (Gaussian $1/\sqrt\pi$) & proved; $\Gamma$-limit cited \cite{esedoglu2017kernels} (flat) & \code{lyapunov\_energy} approximates $\tfrac12 \sum_k \mathrm{Per}(\Omega_k)$ \\
Per-step displacement $\tau\kappa/2$ & proved (leading order); exact disc solution \eqref{eq:disc-exact}; mesh at 0.91--0.97 of it and 0.48--0.54 of the Gaussian, placement-averaged & \S\ref{sec:check} \\
Balanced threshold $\argmax_k(y_{ik} + \psi_k)$ & derived (transportation dual, doc~09); $N{=}2$ case is Ruuth--Wetton & \code{balanced\_assign} $\to$ \code{solve\_dual\_offsets} \\
Dissipation with exact assignment & theorem (cited); the implemented step is the exact MM step for $E^{\Dlump}_\tau$ up to a $\lesssim 10^{-3}$ skew term and the inexact dual & --- \\
Descent inequality \eqref{eq:survives} & proved for any argmax labelling & \code{descent\_slack}; gate G4 \\
Lumped--lumped pairing restores the theorem & stated; \textbf{not implemented} & --- \\
$\tau$ rule \eqref{eq:tau-rule}; freeze \eqref{eq:freeze}; over-merge \eqref{eq:centre} & derived (order of magnitude; $y(0) = 1 - xK_1(x)$ exact for the continuum kernel) & \code{tau\_diagnostics} \\
Freeze ratio \eqref{eq:freeze-ratio} and its two regimes & derived & \code{tau\_diagnostics} \\
Continuation probe & protocol; verdict uncalibrated until NC3-CAL & \code{tau\_continuation\_probe} \\
Init as power diagram of the graph metric & derived & \code{geodesic\_balanced\_init} \\
Ladder transfer error $O(h\,\mathrm{Per})$ & stated & \code{interpolate\_labels} \\
$\Etau$, boundary length, $Q_k$, core loss & instruments & \S\ref{sec:instruments} \\
\bottomrule
\end{tabular}
\end{center}

\section*{Related documents}
\begin{itemize}[itemsep=1pt]
  \item Measurements, pre-registration and the negative controls:\\
    \codefile{docs/experiments/08-mbo-auction-dynamics/} and\\
    \codefile{docs/plans/PHASE1\_BC\_REPLACEMENT\_PLAN.md} (Phase~A).
  \item The transportation dual and the readout it was built for:\\
    \codefile{docs/math/09-balanced-readout/}.
  \item The P1 discretisation and the Phase~1 energy this replaces:\\
    \codefile{docs/math/06-phase1-energy-discretization/}.
  \item Taxonomy of the high-$N$ approaches:\\
    \codefile{docs/reference/PHASE1\_HIGHN\_APPROACHES\_ABCDE.md}.
  \item Held copies of every cited paper: \codefile{docs/papers/README.md}.
\end{itemize}

\bibliographystyle{plain}
\bibliography{../shared/references}

\end{document}
